\documentclass[10pt]{article}
\usepackage[a4paper,margin=1.8cm]{geometry}
\usepackage[T1]{fontenc}
\usepackage{lmodern}
\usepackage{amsmath}
\usepackage{amssymb}
\usepackage{textcomp}
\usepackage{graphicx}
\usepackage{booktabs}
\usepackage{array}
\usepackage{float}
\usepackage{xcolor}
\usepackage{caption}
\usepackage{subcaption}
\usepackage[hidelinks]{hyperref}
\usepackage[normalem]{ulem}
\newcommand{\dead}[1]{\textcolor{gray}{\sout{#1}}}
\setlength{\parindent}{0pt}
\setlength{\parskip}{4pt}
\captionsetup{font=small,labelfont=bf}
\newcommand{\base}{\textbf}
\graphicspath{{run_figures/}{./}}
\newcommand{\CMPIBASE}{0.745}
\newcommand{\CMPIA}{0.748}
\newcommand{\CMPID}{0.750}
\newcommand{\CMPIH}{0.757}
\newcommand{\CMPIO}{0.772}
\newcommand{\CMPIT}{0.773}
\newcommand{\CMPIV}{0.803}
\newcommand{\CMPIFIG}{%
\begin{figure}[H]\centering
\begin{tabular}{@{}c@{\hskip2pt}c@{}}
\includegraphics[width=0.47\textwidth]{Tuning_test_06_Baseline/cmpi.png} & \includegraphics[width=0.47\textwidth]{Tuning_test_06A_fesomA_albpnd028/cmpi.png} \\
\footnotesize Baseline (CMPI 0.745) & \footnotesize 06A (CMPI 0.748) \\
\includegraphics[width=0.47\textwidth]{Tuning_test_06D_HRlike/cmpi.png} & \includegraphics[width=0.47\textwidth]{Tuning_test_06H_fesomH_combo_g_rvice018/cmpi.png} \\
\footnotesize 06D HRlike (CMPI 0.750) & \footnotesize 06H (CMPI 0.757) \\
\includegraphics[width=0.47\textwidth]{Tuning_test_06O_1hcpl_mospp/cmpi.png} & \includegraphics[width=0.47\textwidth]{Tuning_test_06T_1hcpl_mospp_kpplow/cmpi.png} \\
\footnotesize 06O (CMPI 0.772) & \footnotesize 06T (CMPI 0.773) \\
\includegraphics[width=0.47\textwidth]{Tuning_test_06V_1hcpl_mospp_kpplow_entstpc3_1/cmpi.png} \\
\footnotesize 06V (CMPI 0.803) \\
\end{tabular}
\caption{\texttt{part4} CMPI heat maps (per variable/region/season, normalised to CMIP6).}
\end{figure}
}


\title{\vspace{-1.5cm}\textbf{AWI--ESM3 (v3.4, CORE3 / TCO95) PI--control tuning campaign}\\[2pt]
\large \texttt{Tuning\_test\_06*} --- progress report}
\author{Evaluation via \texttt{release\_evaluation\_tool2} (part2 / part8 / part4)}
\date{2026--06--30\ \ (rev.\ 2026--07--06: restructured into three parts.\\
\textbf{Last updated 2026--07--28}: Part~IV added. A feedback-free AMIP reference puts $\approx45\,\%$ of the
TOA imbalance and \emph{all} of the boreal summer cold bias in the atmosphere; CERES shows the summer deficit
is a \emph{land-only cloud} error of $11$\,W/m$^2$ (not the $33$ inferred from CRUNCEP3, which is itself
$\approx21$\,W/m$^2$ too bright), while the global radiation budget already matches CERES to $0.1$\,W/m$^2$.
A controlled two-arm spin-up test shows the forcing change alone --- CRUNCEP\,$\to$\,AMIP, byte-identical
parameters --- costs $54\,\%$ of Siberian tree cover, substantially qualifying the
``competition, not climate'' verdict of \S\ref{sub:decisive}.)}

\begin{document}
\maketitle
\vspace{-0.6cm}

\section*{Update (2026--07--28): new-sea-ice branch \texttt{Tuning\_test\_09A/09B/09C}}

This report now includes the new three-way coupled comparison branch:
\texttt{09A} (baseline + new sea ice), \texttt{09B} (new sea ice + tuned stack), and
\texttt{09C} (new sea ice, 06V variant). All three completed the same 30-year
segment and are compared on the same last-10-year window (model years 1370--1379).

\textbf{Run identities}
\begin{itemize}
\item \texttt{09A}: \texttt{Tuning\_test\_09A\_lpjguess\_Baseline\_coupled\_fromCRUNCEP\_newSeaIce}
\item \texttt{09B}: \texttt{Tuning\_test\_09B\_06T\_1hCPL\_MOSPP\_KPPLOW\_CRUNCEPinit\_newSeaIce}
\item \texttt{09C}: \texttt{Tuning\_test\_09C\_06V\_CRUNCEPinit\_newSeaIce}
\end{itemize}

\textbf{09B tuning block documented on the figures (original $\rightarrow$ tuned)}
\begin{itemize}
\item \texttt{oasis3mct.time\_step}: 7200 s $\rightarrow$ 3600 s
\item \texttt{use\_momix}: \texttt{.false.} $\rightarrow$ \texttt{.true.}
\item \texttt{kpp\_av0}: 0.005 $\rightarrow$ 0.003, \texttt{kpp\_kv0}: 0.005 $\rightarrow$ 0.003
\item \texttt{kpp\_avbckg}: 1e-4 $\rightarrow$ 5e-5, \texttt{kpp\_kvbckg}: 1e-5 $\rightarrow$ 5e-6
\item \texttt{pndaspect}: 0.8 $\rightarrow$ 1.3, \texttt{rfracmax}: 1.0 $\rightarrow$ 0.75, \texttt{albpnd}: 0.2 $\rightarrow$ 0.28
\end{itemize}

\begin{figure}[H]
\centering
\begin{subfigure}[t]{0.49\textwidth}
\centering
\includegraphics[width=\textwidth]{plots/t2m_bias_vs_cruncep3_spinup_09A_09B_09C_nh_last10yrmean.png}
\caption{NH T2m bias vs CRUNCEP3 spinup reference (ocean + land).}
\end{subfigure}\hfill
\begin{subfigure}[t]{0.49\textwidth}
\centering
\includegraphics[width=\textwidth]{plots/treefpc_bias_vs_cruncep3_spinup_09A_09B_09C_nh_last10yrmean.png}
\caption{NH TREEFPC bias vs CRUNCEP3 spinup reference.}
\end{subfigure}
\caption{New-sea-ice branch update: common last-10-year means (1370--1379) for \texttt{09A/09B/09C}.}
\end{figure}

\vspace{-2mm}

\part{Tuning away the high TOA imbalance}
\emph{We set out to cut the AWI--ESM3 PI-control's $+1.5$\,W/m$^2$ TOA imbalance. Every effective lever does it by
\textbf{cooling} the model, and the aggressive coupling/ocean-mixing levers cut the imbalance hardest --- but push the
planet too cold against ERA5. This part documents that trade-off across 22 runs and traces the aggressive lever to a
Southern-Ocean ocean-convection mechanism.}

\section{Goal \& setup}
Reduce the net top-of-atmosphere (TOA) radiative imbalance of the AWI--ESM3 pre-industrial control so it
equilibrates faster, \emph{without} degrading the CMIP6-normalised performance index (CMPI, \texttt{part4})
or inflating the 2\,m-temperature RMSE vs ERA5 (\texttt{part8}). 22 experiments, each a 30-year segment
(model years 1350--1379) of the CORE3/TCO95 configuration (FESOM2 CORE3 mesh, OpenIFS TCO95/A096, LPJ-GUESS),
all branched from one \textbf{Baseline}. The baseline PI-control still drifts with a positive TOA imbalance of
$\approx+1.5$\,W/m$^2$ (energy \emph{into} the system --- the signature of unequilibrated spin-up and/or ongoing
ocean heat uptake). \textbf{A caution on the target.} The tuning levers change the model's \emph{equilibrium}, not
its rate of approach, so ``lower imbalance'' is a useful \emph{ranking} objective but not automatically ``more
equilibrated'': a low-imbalance-but-cold state is not obviously better than a high-imbalance-but-unbiased one. This
is why the campaign's real aim (\S\ref{sec:rec}) is a \emph{warm-floor} state where the imbalance is cut by a
\emph{moderate} lever, not by cooling into a cold bias --- i.e.\ we treat imbalance and 2\,m-T bias as joint targets,
not imbalance alone.

\textbf{Structural note.} Almost every non-baseline run also carries a \emph{common melt-pond shift}
(\texttt{albpnd}\,$0.2\!\to\!0.28$, \texttt{rfracmax}\,$1.0\!\to\!0.75$, \texttt{pndaspect}\,$0.8\!\to\!1.3$).
Differences below are each run's \emph{distinguishing} change on top of that.

\textbf{Three parts.} \textbf{(I)} evaluates the 22 imbalance-tuning runs and exposes the imbalance$\leftrightarrow$cold-bias
trade-off, tracing the aggressive lever to a Southern-Ocean mechanism. \textbf{(II)} shows that the cold bias which
makes that lever costly is largely \emph{not} the tuning but a \textbf{structural} LPJ-GUESS boreal-forest collapse
(\S\ref{sec:lpjg}) --- a \emph{counter-lever}: fixing it warms the NH high latitudes, buying the RMSE headroom to spend
the aggressive global cooling on the imbalance. A decisive feedback-free experiment pins that fix to vegetation, not
climate. \textbf{(III)} combines them into a work plan: raise the structural warm floor, then cut the imbalance
(\S\ref{sec:rec}).

\section{Headline: the imbalance\,$\leftrightarrow$\,cold-bias trade-off}
Every lever that meaningfully lowers the TOA imbalance does so by \emph{cooling} the model, which worsens the
cold bias against ERA5. Two families separate cleanly: the sea-ice/albedo/melt-pond changes move gently
down-left (small imbalance gain, small bias cost); the coupling + ocean-mixing changes cut the imbalance hard
(30--57\,\%) but push the planet $1.7$--$2.1$\,K too cold.

\textbf{The cold-bias axis is two things at once.} It mixes (i) a \emph{shared structural floor} --- every run,
including the baseline and the ``clean'' melt-pond runs, carries a strong Siberia/N.\ Canada cold bias that is
\emph{not} a tuning artefact but an LPJ-GUESS boreal-forest-cover bug (\S\ref{sec:lpjg}); and (ii) the \emph{differential},
tuning-induced global cooling that the aggressive runs add on top. The structural floor is the same for all
points and is independently fixable; the tuning only moves runs vertically relative to it. This is why the
``warming lever'' in \S\ref{sec:rec} is the LPJG fix, not another atmospheric retune.

\begin{figure}[H]
\centering
\includegraphics[width=0.84\textwidth]{plots/summary_scatter_official.png}
\caption{Net TOA imbalance (30-yr mean, \texttt{part2}) vs 2\,m-T RMSD against ERA5 (\texttt{part8}).
Dotted lines = baseline; lower-left is better. Colour = parameter family.}
\end{figure}

\begin{figure}[H]
\centering
\includegraphics[width=0.92\textwidth]{plots/summary_TOA_timeseries.png}
\caption{Annual global-mean net TOA imbalance, all 22 runs (independent CDO calculation; matches \texttt{part2}).
None has flattened within 30\,yr, so treat the mean level as a \emph{ranking}, not a final equilibrium value.}
\end{figure}

\section{Master metric table (all 22 runs, sorted by TOA imbalance)}
\begin{center}\small
\begin{tabular}{l r r r r r}
\toprule
Run & TOA imb.\ [W/m$^2$] & vs base & 2m RMSD [K] & 2m bias [K] & CMPI \\
\midrule
06V\_1hcpl\_mospp\_kpplow\_entstpc3\_1 & 0.655 & $-57\%$ & 2.671 & $-2.124$ & 0.803 \\
06T\_1hcpl\_mospp\_kpplow             & 0.952 & $-38\%$ & 2.385 & $-1.791$ & 0.773 \\
06O5\_1hcpl\_mospp\_kv002             & 0.968 & $-37\%$ & 2.511 & $-1.904$ & 0.788 \\
06O\_1hcpl\_mospp                     & 0.978 & $-36\%$ & 2.389 & $-1.774$ & 0.772 \\
06O4\_1hcpl\_mospp\_kv0012            & 1.007 & $-34\%$ & 2.387 & $-1.807$ & 0.781 \\
06U\_1hcpl\_mospp\_kpplow\_openw0075  & 1.020 & $-33\%$ & 2.422 & $-1.845$ & 0.776 \\
06M\_1hcpl                            & 1.062 & $-31\%$ & 2.315 & $-1.658$ & 0.773 \\
06P\_entstpc3\_1                      & 1.091 & $-29\%$ & 2.076 & $-1.358$ & 0.759 \\
06H\_combo\_g\_rvice018               & 1.220 & $-20\%$ & 2.082 & $-1.394$ & 0.757 \\
06G\_stronger\_meltpond               & 1.274 & $-17\%$ & 1.990 & $-1.274$ & 0.753 \\
06L\_extra\_strong\_meltpond          & 1.320 & $-14\%$ & 2.050 & $-1.332$ & 0.751 \\
06E\_rvice018\_ggaussb\_m06           & 1.335 & $-13\%$ & 1.969 & $-1.231$ & 0.754 \\
06D\_HRlike                           & 1.377 & $-10\%$ & 1.892 & $-1.146$ & 0.750 \\
06F\_lrdalb\_true                     & 1.378 & $-10\%$ & 1.962 & $-1.219$ & 0.752 \\
06N\_mospp                            & 1.382 & $-10\%$ & 2.011 & $-1.310$ & 0.751 \\
06C\_albpnd028\_rfrac075\_pndaspect13 & 1.391 & $-9\%$  & 1.884 & $-1.109$ & 0.751 \\
06Q\_ralbsead0045                     & 1.397 & $-9\%$  & 1.891 & $-1.154$ & 0.748 \\
06S\_evp0                             & 1.419 & $-7\%$  & 1.935 & $-1.191$ & 0.745 \\
06R\_openwater\_albedo0075            & 1.443 & $-6\%$  & 1.860 & $-1.127$ & 0.744 \\
06B\_albpnd028\_rfrac075              & 1.447 & $-5\%$  & 1.847 & $-1.062$ & 0.741 \\
06A\_albpnd028                        & 1.457 & $-5\%$  & 1.811 & $-1.004$ & 0.748 \\
\textbf{06\_Baseline}                 & \textbf{1.531} & -- & \textbf{1.717} & \textbf{$-0.939$} & \textbf{\CMPIBASE} \\
\bottomrule
\end{tabular}
\end{center}
{\footnotesize TOA = 30-yr global-mean net TOA from \texttt{part2\_rad\_balance.py}. RMSD/bias from \texttt{part8}
methodology (cos-lat weighted, model$-$ERA5, yrs 1354--1379); cross-checked vs the tool's on-figure values
(baseline 1.717 vs 1.721; 06V 2.671 vs 2.678). CMPI from \texttt{part4\_cmpi.py}: $<1$ = better than CMIP6
ensemble-mean error; lower is better.}

The per-run distinguishing parameter change (which knob each \texttt{06*} run alters on top of the common melt-pond
shift) is catalogued in App.~\ref{app:params}.

\clearpage
\section{Core diagnostics: radiation, 2\,m temperature, CMPI}

\subsection{Radiation budget (\texttt{part2})}\label{sec:rad}
Surface net flux tracks TOA to within $\approx0.1$\,W/m$^2$ in every run --- i.e.\ the atmospheric column stores
essentially no energy on these multi-year means, as expected. Note this does \emph{not} by itself localise the
\emph{cause}: a TOA\,$\approx$\,surface match is equally consistent with a genuine TOA radiative error and with
ocean heat uptake / unequilibrated spin-up drift passing cleanly through the surface (\S\ref{sec:so},
\S\ref{sub:withwithout}). What it rules out is an atmosphere-only bookkeeping leak, so the imbalance and the surface
energy budget can be read interchangeably. Baseline $\approx+1.53$\,W/m$^2$; the coupling+mixing stack drives it
down to $+0.66$\,W/m$^2$ (06V).

\begin{figure}[H]
\centering
\begin{tabular}{@{}c@{\hskip2pt}c@{}}
\includegraphics[width=0.47\textwidth]{Tuning_test_06_Baseline/radiation_budget.png} &
\includegraphics[width=0.47\textwidth]{Tuning_test_06A_fesomA_albpnd028/radiation_budget.png} \\
\footnotesize Baseline (TOA 1.53) & \footnotesize 06A albpnd0.28 (1.46) \\
\includegraphics[width=0.47\textwidth]{Tuning_test_06H_fesomH_combo_g_rvice018/radiation_budget.png} &
\includegraphics[width=0.47\textwidth]{Tuning_test_06P_fesomP_entstpc3_1/radiation_budget.png} \\
\footnotesize 06H meltpond+rvice (1.22) & \footnotesize 06P entstpc3 (1.09) \\
\includegraphics[width=0.47\textwidth]{Tuning_test_06O_1hcpl_mospp/radiation_budget.png} &
\includegraphics[width=0.47\textwidth]{Tuning_test_06V_1hcpl_mospp_kpplow_entstpc3_1/radiation_budget.png} \\
\footnotesize 06O 1hcpl+mospp (0.98) & \footnotesize 06V (most aggressive, 0.66) \\
\end{tabular}
\caption{\texttt{part2} radiation budget: grey = net TOA, blue = net surface, dark = their difference (W/m$^2$).}
\end{figure}

\clearpage
\subsection{2\,m temperature bias vs ERA5 (\texttt{part8})}
Baseline: cold over NH continents, warm over the Southern Ocean and tropical upwelling. The aggressive coupling
runs remove the Southern-Ocean warm bias but turn the whole planet cold (06V is cold almost everywhere).
\emph{Caveat:} this is PI-control vs present-day ERA5, so $\approx0.5$--$0.7$\,K of global cold offset is
physically expected; the coupling runs' $-1.7$ to $-2.1$\,K is clearly too much, the melt-pond family
($-1.0$ to $-1.4$\,K) is more defensible.

The deep \textbf{Siberia / N.\ Canada cold lobe is present in \emph{every} panel} --- baseline and the clean
melt-pond runs included --- and barely changes between them: it is the structural LPJ-GUESS forest-cover signal
(\S\ref{sec:lpjg}), not a product of the imbalance tuning. The aggressive coupling runs add a broad, near-global cooling
\emph{on top} of that fixed lobe.

\begin{figure}[H]
\centering
\begin{tabular}{@{}c@{\hskip2pt}c@{}}
\includegraphics[width=0.47\textwidth]{Tuning_test_06_Baseline/t2m_vs_ERA5.png} &
\includegraphics[width=0.47\textwidth]{Tuning_test_06A_fesomA_albpnd028/t2m_vs_ERA5.png} \\
\footnotesize Baseline (RMSD 1.72, bias $-0.94$) & \footnotesize 06A (1.81, $-1.00$) \\
\includegraphics[width=0.47\textwidth]{Tuning_test_06H_fesomH_combo_g_rvice018/t2m_vs_ERA5.png} &
\includegraphics[width=0.47\textwidth]{Tuning_test_06D_HRlike/t2m_vs_ERA5.png} \\
\footnotesize 06H (2.08, $-1.39$) & \footnotesize 06D HRlike (1.89, $-1.15$) \\
\includegraphics[width=0.47\textwidth]{Tuning_test_06O_1hcpl_mospp/t2m_vs_ERA5.png} &
\includegraphics[width=0.47\textwidth]{Tuning_test_06V_1hcpl_mospp_kpplow_entstpc3_1/t2m_vs_ERA5.png} \\
\footnotesize 06O (2.39, $-1.77$) & \footnotesize 06V (2.67, $-2.12$) \\
\end{tabular}
\caption{\texttt{part8} 2\,m air-temperature bias (model $-$ ERA5), EqualEarth, $\pm8$\,K scale.}
\end{figure}

\clearpage
\subsection{CMPI (\texttt{part4})}\label{sec:cmpi}
CMPI integrates 13 fields (siconc, tas, clt, pr, rlut, uas, vas, ua@300, zg@500, zos, mlotst, thetao, so) over
regions/seasons, normalised against 30 CMIP6 models; \textbf{$<1$ = better than the CMIP6 ensemble-mean error}.
Run on a subset spanning the trade-off front (now complete for all 22 runs).

\textbf{How to read this (important).} CMPI here compares a \emph{PI-control} run to \emph{present-day}
observations, so a moderate global cold offset --- and hence some CMPI rise --- is \emph{expected} and is not
by itself a disqualifier. The informative question is not ``did CMPI go up?'' but ``\emph{where} did the error
change, and is that consistent with the parameter that was altered?'' A sea-ice-albedo change should move
\texttt{siconc} in the polar regions and little else; if it instead degrades the tropics, that is a red flag.
The per-variable / per-region breakdown below is the tool for that judgement; the global number is a summary,
not a verdict.

\begin{center}\small
\begin{tabular}{l r r r}
\toprule
Run & CMPI & TOA imb. & 2m RMSD \\
\midrule
06\_Baseline          & 0.745 & 1.531 & 1.717 \\
06A\_albpnd028        & 0.748 & 1.457 & 1.811 \\
06D\_HRlike           & 0.750 & 1.377 & 1.892 \\
06H\_combo\_g\_rvice018 & 0.757 & 1.220 & 2.082 \\
06O\_1hcpl\_mospp     & 0.772 & 0.978 & 2.389 \\
06T\_1hcpl\_mospp\_kpplow & 0.773 & 0.952 & 2.385 \\
06V\_..\_entstpc3\_1  & 0.803 & 0.655 & 2.671 \\
\bottomrule
\end{tabular}
\end{center}

\CMPIFIG

\subsubsection{Where does the error change? (per-variable / per-region)}
The maps below show each run's CMPI fractional error \emph{minus baseline} (red = worse, blue = better).
This is where the PI-vs-PD caveat earns its keep: the global CMPI rises monotonically
(0.745\,$\to$\,0.803), but the \emph{pattern} of change tells a much more favourable story for the targeted
tunings.

\begin{figure}[H]
\centering
\includegraphics[width=0.86\textwidth]{plots/cmpi_breakdown_by_variable.png}\\[2pt]
\includegraphics[width=0.74\textwidth]{plots/cmpi_breakdown_by_region.png}
\caption{Change in CMPI fractional error vs Baseline, by variable (top) and region (bottom).}
\end{figure}

\textbf{Reading the breakdown:}
\begin{itemize}\itemsep1pt
\item \textbf{06A / 06D (melt-pond):} improve \texttt{siconc} ($-0.02$/$-0.08$) with essentially \emph{zero}
collateral elsewhere and a flat region map --- a cleanly targeted change. These should be regarded as
``no degradation where it matters,'' despite the tiny global CMPI uptick.
\item \textbf{06O / 06T (mospp\,+\,kpplow):} genuinely \emph{improve} the ocean interior (\texttt{thetao},
\texttt{so}: $\approx-0.04$/$-0.05$), \texttt{siconc} and \texttt{clt} --- exactly where ocean-mixing should act.
But they \emph{degrade} \texttt{tas}/\texttt{zg}/\texttt{rlut} ($+0.16$ to $+0.26$) and even their own target
\texttt{mlotst} ($+0.17$), and the region map reddens \emph{everywhere incl.\ tropics and Ni\~no-3.4} --- i.e.\
the cooling leaks well beyond where an ocean-mixing parameter has any business acting. \emph{That} is the
warning sign, more than the global number.
\item \textbf{06V (+entstpc3):} largest, most widespread atmospheric degradation (\texttt{tas} $+0.45$,
\texttt{zg} $+0.46$); the deep-convection change destabilises the whole atmospheric column.
\end{itemize}

\clearpage
\section{Southern-Ocean mechanism (adapted from a270234's notebooks)}\label{sec:so}
To see \emph{why} the ocean-mixing/coupling levers act as they do, two of a270234's Weddell-Sea diagnostics
(\texttt{mld\_sic\_sit\_compare}, \texttt{temp\_fesom\_core3\_WWS}) were re-pointed at the coupled
\texttt{Tuning\_test\_06} runs (same CORE3 mesh). Panels show the Baseline absolute field plus each run
\emph{minus} Baseline, austral winter (September), yr1370--79. (The \texttt{temp\_fesom\_core3\_WWS} tripyview
transect could not decode these runs' pre-1582 \texttt{gregorian} calendar, so an equivalent pyfesom2
depth--latitude section was used.)

\begin{figure}[H]\centering
\includegraphics[width=0.82\textwidth]{plots/mld_winter_antarctic.png}
\caption{Antarctic September mixed-layer depth (MLD3): Baseline \emph{absolute} (top-left, own scale), others $=$
\emph{anomaly} vs Baseline on a symmetric $\pm70$\,m diverging scale (98th-pctile auto-scale; the intense
Weddell/Ross convection core saturates it). The mixing runs (06N/O/T/V) restructure the Weddell/Ross winter mixing
relative to Baseline, whereas melt-pond (06A) barely changes it. Deep open-ocean winter mixing there is FESOM
``open-ocean convection'' that ventilates warm deep water --- a known driver of the Southern-Ocean warm bias
(Heuz\'e et al.\ 2013; Heuz\'e 2021) and of TOA imbalance.}
\end{figure}

\begin{figure}[H]\centering
\includegraphics[width=0.82\textwidth]{plots/temp_wws_section_tuning.png}
\caption{West Weddell Sea (50--30$^{\circ}$W) September potential-temperature section: Baseline absolute and
anomalies. The ventilated mid-column cools in the mixing runs (less convective entrainment of warm deep water to the
surface), consistent with suppressed open-ocean convection; note that suppressing convection should leave the
\emph{deepest} water trapped/warmer rather than colder, and a near-surface edge warm anomaly is also present ---
i.e.\ the signal is a redistribution between layers, not uniform subsurface cooling.}
\end{figure}

\begin{figure}[H]\centering
\includegraphics[width=0.92\textwidth]{plots/sit_winter_antarctic.png}
\caption{Antarctic September sea-ice \emph{thickness} (Baseline absolute, top-left; others $=$ anomaly vs Baseline).
Thickness --- not concentration --- is the honest growth diagnostic here: winter concentration is already near
$1$ over most of the pack, so it \emph{saturates} and cannot register added ice, whereas thickness is unbounded. The
mixing runs (06O/06T/06V) thicken the pack broadly (positive/red anomaly over the Weddell and Ross sectors),
consistent with suppressed open-ocean convection insulating and stratifying the surface; melt-pond (06A) barely
changes it. (The concentration version, \texttt{sic\_winter\_antarctic.png}, shows only the ice-edge shifts for the
reason above.)}
\end{figure}

\textbf{Mechanism.} The mospp / weaker-KPP changes reduce Weddell-Sea open-ocean winter convection, cooling the
subsurface and letting sea ice grow --- improving \texttt{thetao}/\texttt{so}/\texttt{siconc} (consistent with the
CMPI breakdown) and removing the Southern-Ocean surface warm bias. Spurious open-ocean deep convection in the
Weddell/Ross gyres is the canonical CMIP Antarctic bias (Heuz\'e et al.\ 2013; Heuz\'e 2021, which evaluates AWI--CM
in the multi-model set; Mohrmann et al.\ 2021 for the polynya surface expression), and open-ocean Weddell convection
is a real-but-rare mode that models over-express (de Lavergne et al.\ 2014) --- so damping it is physically the right
direction. The \emph{proximate radiative} closure is defensible: more winter ice $\to$ higher albedo $\to$ less
absorbed SW $\to$ lower imbalance. We are more cautious about the ``reduced ocean heat uptake'' half of the chain:
suppressing convection also \emph{traps} heat at depth, which can \emph{increase} deep ocean-heat-content drift, so
the net uptake change is not established here --- it should be pinned to an OHC-tendency diagnostic before being
asserted. The cost is that the surface cooling propagates globally into the atmosphere (\texttt{tas}/\texttt{zg}),
the broad cold bias seen in \texttt{part8}.

\textbf{Counter-literature (a compensating-error risk this campaign must own).} \S\ref{sec:so} attributes the SO warm
bias and its removal almost entirely to \emph{ocean} convection, but Hyder et al.\ (2018, \emph{Nat.\ Commun.})
trace the leading $40$--$60$\textdegree S warm-SST bias primarily to \emph{atmospheric} net surface-heat-flux errors
dominated by shortwave/cloud biases, and warn that a small \emph{net} flux bias is often achieved by
\emph{compensating} component errors. The two coexist at different scales --- broad atmospheric-SW warmth vs.\
localised Weddell/Ross deep-convection anomalies --- and our mixing levers act only on the ocean branch. If part of
the bias is atmospheric-SW in origin, the ocean-mixing lever is partly compensating for an atmospheric error, which
is \emph{exactly} the ``cooling leaks beyond where the parameter should act'' failure mode flagged for these runs in
\S\ref{sec:cmpi} (the region map reddens into the tropics/Ni\~no-3.4). This tension is the report's own evidence
pointing at Hyder; it strengthens, not weakens, the case for treating aggressive ocean-mixing tunings as
provisional.

\clearpage
\part{Counter-tuning the NH cold bias --- the structural floor}
\emph{The cold bias that makes Part~I's aggressive lever expensive is largely \textbf{not} the tuning: it is a
structural LPJ-GUESS boreal-forest collapse that every run inherits. That makes it a \textbf{counter-lever} --- fixing
it warms the NH high latitudes at essentially unchanged imbalance, buying the RMSE headroom to spend the aggressive
global cooling on the imbalance. This part attributes the bias with a clean WITH-vs-WITHOUT-LPJG experiment, resolves
its two mechanisms and a residual, and --- with a decisive feedback-free spin-up --- pins the fix to vegetation
competition, not climate.}

\section{The NH high-latitude cold bias is a structural LPJ-GUESS problem}\label{sec:lpjg}
The runs carry a severe near-surface \textbf{cold bias over Siberia and northern Canada} (visible in every
\texttt{part8} map). This chapter documents a working investigation into whether the coupled dynamic vegetation
(LPJ-GUESS) is responsible --- motivated by the idea that fixing it could be a \emph{warming} lever to offset the
imbalance-tuning cold bias, and by prior EC-Earth experience that LPJG drives such boreal biases.

\subsection{Symptom: boreal needleleaf replaced by grass}
Dominant-PFT maps (tool \texttt{part26}) for the HTESSEL spin-up (the reference, run with the same atmosphere but
\emph{without} atmosphere$\leftrightarrow$LPJG feedback) vs the coupled runs show boreal needleleaf giving way to
C3 grass over Siberia / N.\ Canada. Tree-dominant fraction north of 50\textdegree N: \textbf{20.1\,\% (spin-up)
$\to$ 16.8\,\% (baseline) $\to$ 11.4\,\% (06V)} --- progressively worse with the more aggressive (colder) runs.
\emph{Caveat on the prior:} DGVMs disagree substantially on boreal tree-vs-grass partitioning, and real boreal
forest is genuinely open with substantial non-tree ground cover (Jiang et al.\ 2012), so ``grass over Siberia'' is
not \emph{automatically} a bug. What makes the establishment-failure reading defensible here is the \emph{magnitude}:
the coupled state reaches cvh\,$\approx$\,0.02--0.05 with \dead{no closed canopy anywhere (FORESTFPC$=0$)}\footnote{\textbf{Void diagnostic.} \texttt{FORESTFPC} sums only PFTs with \texttt{landcover==FOREST} (\texttt{commonoutput.cpp}:1143), i.e.\ the \emph{managed-forest} land-use class; natural trees are counted in \texttt{TREEFPC} via the \texttt{NATURAL} branch. In this PI configuration nothing is allocated to FOREST (\texttt{Natural\_sum}$=0.684$, \texttt{Total}$=0.665$), so \texttt{FORESTFPC}$\equiv0$ \emph{by construction} while \texttt{TREEFPC}$=0.323$. It is a land-use bookkeeping artifact and says nothing about canopy closure.} (
\S\ref{sec:workplan}), an outlier well below the $\sim0.4$--$0.7$ closed-taiga range (MODIS VCF), not merely a
grass-leaning partition.

\begin{figure}[H]\centering
\includegraphics[width=0.98\textwidth]{plots/pft_boreal_compare.png}
\caption{Dominant PFT, boreal belt: HTESSEL spin-up (reference) vs coupled baseline vs 06V. Needleleaf (blues)
thins and C3 grass (yellow) expands across central/eastern Siberia in the coupled runs.}
\end{figure}

\subsection{Is it grass growing, or trees dying? And is it a year-1 jump or a slow drift?}
Belt-mean (55--75\textdegree N) needleleaf LAI is essentially \emph{conserved} over the baseline run (0.40$\to$0.37,
trend $\approx$0), while C3 grass LAI \emph{rises} (0.65$\to$0.83, $+$0.10/dec): the signal is \textbf{grass
growing}, with needleleaf loss only regional (E.\ Siberia) and offset by gains in the milder west. The coupled
run is initialised directly from the spin-up's year-3850 LPJG state, and baseline year-1350 boreal LAI
(tree 0.407, grass 0.649) is \emph{identical} to the spin-up end-state (0.401 / 0.669) --- so there is
\textbf{no wholesale year-1 collapse}; the degraded, grass-heavy state is \emph{inherited from the spin-up}.
Within 1350--1379 the boreal 2\,m temperature is at its (cold) level from year one and does \emph{not} drift
down (trend $+$0.42\,K/dec, noisy) --- consistent with a memory-less atmosphere that equilibrates to the surface
state within a year rather than a slow vegetation runaway. (These are belt-mean \emph{LAI} diagnostics; as
\S\ref{sub:withwithout} shows, the quantity that actually matters is not LAI-per-tile but the forest \emph{cover fraction}.)

\subsection{The decisive test: a clean WITH-vs-WITHOUT-LPJG pair (CORE3, 30\,yr)}\label{sub:withwithout}
The earlier control was confounded by the ocean (core2 vs core3). We removed that confound: a purpose-built
\textbf{no-LPJG AWI-CM3-develop reference} --- \emph{same} CORE3 ocean, \emph{same} TCO95 atmosphere and matched
component branches, HTESSEL\,+\,prescribed satellite vegetation, a full 30-yr PI-control (model years 1350--1379,
\texttt{awicm3\_noLPJG\_CORE3\_30y}) --- differs from the LPJG Baseline in \textbf{LPJG alone}, so the 2\,m-T
difference is now a clean attribution. Its net TOA imbalance is $+1.73$\,W/m$^2$ (\texttt{part2}), near the
warm/high-imbalance end of Fig.~1 --- confirming the LPJG signal is a near-surface \emph{land} effect, orthogonal
to the TOA tuning of \S2. \textbf{Crucially, it beats the LPJG Baseline on \emph{both} skill metrics: global 2\,m-T RMSD vs ERA5 1.514\,K
(bias $-0.76$\,K, \texttt{part8}) vs 1.717/$-0.94$, and CMPI 0.738 (\texttt{part4}) vs 0.745} --- at higher
(warmer) imbalance. Removing LPJG improves the global temperature fit \emph{and} the multivariate skill score
\emph{while} the imbalance changes by only $+0.2$\,W/m$^2$ ($+1.73$ vs $+1.53$) --- within the inter-run drift of
these unequilibrated 30-yr segments (\S\ref{sec:caveats}), so the warming lever is real \emph{and costs essentially
no imbalance to within that uncertainty} (not, strictly, ``zero''). Over the final decade, boreal DJF:
\begin{itemize}\itemsep1pt
\item \textbf{Forest cover collapses.} Siberian high-veg cover \texttt{cvh}: \textbf{0.63 (no LPJG) $\to$ 0.02
(LPJG)} (pan-boreal 0.49\,$\to$\,0.07). LPJG passes essentially no forest to the atmosphere.
\item \textbf{And it is genuinely colder --- confound gone.} With the ocean held fixed, the LPJG run is
unambiguously \textbf{$-3.1$\,K} colder over Siberia in DJF ($-30.4$ vs $-27.4$\textdegree C).
\item \textbf{Snow-masking (spring albedo)} operates as before: forest removal brightens the snow season
(MAM \textbf{+0.077}). This is the canonical boreal-forest climate mechanism --- a masked snowpack lowers albedo and
warms high latitudes, and its removal cools most in spring (Bonan, Pollard \& Thompson 1992; Betts 2000; Bonan 2008),
with tree/canopy cover setting the snow-albedo-feedback magnitude that many models misrepresent (Loranty et al.\
2014; Chapin et al.\ 2005) --- so a model that loses its forest gets exactly this feedback wrong. But it is a
\emph{sunlight} effect and cannot explain the equally-cold polar night (\S\ref{sub:seasonality}).
\end{itemize}

\begin{figure}[H]\centering
\includegraphics[width=0.99\textwidth]{plots/djf_roughness_inversion_maps.png}
\caption{DJF WITH$-$WITHOUT LPJG, boreal NH (final decade, land-masked). Forest-cover collapse (\texttt{cvh}, top-left,
brown) is spatially co-located with the cold anomaly (top-mid), a \emph{stronger} low-level inversion
(T$_{925}-$T$_{2m}$, top-right) and surface decoupling (T$_{2m}-$T$_{skt}$, bottom-left), at \emph{unchanged} cloud
(bottom-mid). LW-down (bottom-right) is modestly reduced over the cold interior.}
\end{figure}

\subsection{Seasonality: which months carry the cold effect}\label{sub:seasonality}
Resolving the same clean CORE3 WITH$-$WITHOUT pair by \emph{month} (boreal box 55--70\textdegree N, Siberia + N.\ America,
final decade) shows the LPJG cold effect is \textbf{present in every month except high summer}. The WITH$-$WITHOUT
2\,m-T difference is negative from September through June and only vanishes in July--August
($-0.1$/$-0.2$\,K), peaking in the \emph{shoulder/dark} seasons: \textbf{November $-3.9$\,K, April $-3.3$\,K},
with DJF $-2.4$\,K and MAM $-2.5$\,K (annual box-mean $-2.0$\,K). \textbf{This is the key seasonal fingerprint}: a
pure snow-albedo (sunlit) mechanism would confine the cold to the bright spring, yet the effect is
\emph{as large in near-sunless November/December} --- which is exactly what motivates the winter
roughness$\to$inversion mechanism of \S\ref{sub:roughness}. The near-zero July--August signal also confirms the
growing-season atmosphere is unaffected (consistent with JJA gates being cleared, \S\ref{sub:withwithout}).
Overlaid is the independent \textbf{OIFS--AMIP} run (\S\ref{sub:amipgen}: observed-SST, no-LPJG, PI 1870--79) --- it
tracks the \emph{warm} WITHOUT curve (annual $-7.6$ vs $-6.3$\,\textdegree C; DJF $-25.6$ vs $-24.9$; MAM even milder),
confirming that a feedback-free, realistically-forced boreal climate is warm, and that the forcing we drive the
decisive offline re-spin with (\S\ref{sub:amipgen}) is \emph{not} itself cold-biased.

\begin{figure}[H]\centering
\includegraphics[width=0.99\textwidth]{plots/nolpjg_seasonal_cycle.png}
\caption{Boreal (55--70\textdegree N, Siberia+Canada) monthly climatology. \emph{Left:} 2\,m-T annual cycle ---
WITHOUT LPJG (green, clean CORE3 reference), WITH LPJG (red, coupled), and the OIFS--AMIP obs-SST no-LPJG forcing
run (blue); blue bars = WITH$-$WITHOUT [K, right axis], negative (LPJG colder) in every month bar July--August,
strongest in Nov/Apr. \emph{Right:} surface albedo ($1-$SSR/SSRD); LPJG brightens the snow season (spring/autumn),
the sunlit part of the bias. (AMIP emits no down-welling SW, so it appears on the temperature panel only. WITH/WITHOUT
are model years 1370--79; AMIP is obs-SST 1870--79 --- both $\sim$PI.)}
\end{figure}

\subsection{A second, winter mechanism: roughness $\to$ inversion}\label{sub:roughness}
Snow-albedo needs sunlight, yet the cold bias is as strong in DJF as in MAM. The missing forest acts a second way,
in the dark. OpenIFS roughness length is cover-weighted by \texttt{cvh}, and needleleaf forest ($z_{0m}=2.0$\,m) is
$20$--$60\times$ rougher than the tundra/grass ($0.034$--$0.10$\,m, \texttt{susveg\_mod.F90}) that replaces it. A
smoother surface couples less efficiently to the atmosphere; in polar night the stable boundary layer (SBL) decouples
and the surface radiatively cools, deepening the near-surface inversion. That surface--atmosphere coupling strength in
the SBL controls screen temperature and inversion depth is well established (Holtslag et al.\ 2013; Sandu et al.\
2013); we assemble the specific vegetation-roughness\,$\to$\,SBL-decoupling\,$\to$\,screen-T chain from the
land-surface and SBL literatures rather than a single published result, so it is presented as a mechanism the data
then test, not a settled citation. \textbf{The data show exactly this fingerprint}
(DJF Siberia, LPJG\,$-$\,no-LPJG): a \textbf{+2.1\,K stronger} low-level inversion (T$_{925}-$T$_{2m}$), a
\textbf{+0.7\,K stronger} surface inversion (T$_{2m}-$T$_{skt}$), \textbf{unchanged cloud} ($\Delta\approx0.00$) and
\textbf{$-4.7$\,W/m$^2$ less} down-welling longwave --- a decoupled, radiatively-cooling stable surface, not a
cloud/radiation artefact. Spatially over boreal land the forest-loss pattern \emph{is} the inversion pattern:
area-weighted $r(\Delta\texttt{cvh},\Delta\text{inv})=-0.60$, $r(\Delta\texttt{cvh},\Delta\text{sfc-decouple})=-0.68$,
$r(\Delta T_{2m},\Delta\text{inv})=-0.83$, while the negative control $r(\Delta\texttt{cvh},\Delta\text{cloud})=-0.05$
is $\approx0$ --- a real mechanism, not a shared trend. Regression lever: \textbf{restoring \texttt{cvh} by $+0.4$
warms DJF screen-T by $\sim+1.2$\,K} through roughness alone and weakens the inversion by $\sim0.7$\,K --- additive
to the spring/albedo warming. This slope-based estimate inherits the scatter of the underlying correlation
($r=0.63$, so $<40\%$ of grid-cell variance explained); read it as $\sim+1.2\pm0.5$\,K, an order-of-magnitude lever
rather than a precise prediction. The \emph{direction} and existence of the effect rest on the WITH$-$WITHOUT
experimental design (\S\ref{sub:withwithout}), not on this regression.

\begin{figure}[H]\centering
\includegraphics[width=0.99\textwidth]{plots/djf_roughness_scatter.png}
\caption{DJF boreal-land, WITH$-$WITHOUT LPJG (final decade). Forest loss (negative $\Delta$\texttt{cvh}) tracks the
cold ($\Delta T_{2m}$, left), the stronger inversion (middle) and, transitively, cold\,$\leftrightarrow$\,inversion
(right, $r=-0.83$). Colour = cos-lat weight.}
\end{figure}

\subsection{How much does the LPJG fix buy? (vs ERA5)}\label{sub:levervsera5}
Against ERA5 the current Baseline is \textbf{$-4.5$\,K} too cold over Siberia in DJF. The lever helps but does not
close it: a realistic partial fix ($+1$\,K) leaves $-3.5$\,K, and \emph{full} forest restoration (the no-LPJG run
itself --- the ceiling) reaches \textbf{$-1.4$\,K}. So vegetation/roughness accounts for $\sim$3\,K, roughly
two-thirds of the Siberian DJF cold bias, and boreal-land RMSD falls 4.4\,$\to$\,3.3\,K --- but a residual $\sim$1.5\,K
survives \emph{even with perfect forest}.

\begin{figure}[H]\centering
\includegraphics[width=0.99\textwidth]{plots/djf_bias_vs_era5.png}
\caption{DJF 2\,m-T bias vs ERA5, boreal NH (final decade): current LPJG Baseline (left, model output), an
\textbf{illustrative} $+1$\,K roughness lever (middle --- a uniform $+1$\,K offset applied to the Baseline, \emph{not}
a simulation), and full forest restoration = the no-LPJG run (right, model output). The lever panel only visualises
the magnitude of a plausible partial fix; the load-bearing comparison is left vs right, which shows a residual cold
survives even full forest restoration.}
\end{figure}

\subsection{The residual driver: an over-strong stable-boundary-layer inversion}\label{sub:residual}
Diagnosing that residual in the \emph{full-forest} run (vegetation thereby ruled out) against ERA5 (DJFM, boreal
land) points to a single process. The model's low-level inversion is \textbf{+1.4\,K too strong}
(T$_{925}-$T$_{2m}$: model $+5.6$ vs ERA5 $+4.2$\,K), and the excess is \textbf{entirely surface-trapped} --- the
$925$--$1000$\,hPa inversion bias is $+0.07$\,K ($\approx0$), so all the extra cooling sits in the lowest layer.
Spatially this inversion-excess explains \textbf{$r=+0.76$} of \emph{where} the model is too cold; the Siberian
High (MSLP bias $+0.15$\,hPa, $r=0.27$) and cloud-\emph{fraction} (bias $+0.07$, $r=0.00$) are ruled out. A
too-strong wintertime polar inversion is a real, observationally-benchmarked model bias (Medeiros et al.\ 2011;
Zhang et al.\ 2011), which supports reading the ``$+1.4$\,K too strong'' residual as genuine. That surface-trapped
signature, with circulation and cloud-fraction eliminated, is also the textbook fingerprint of a \textbf{stable
boundary layer that under-mixes} --- the IFS-family ``too-stable'' winter cold bias (Sandu et al.\ 2013). Its lever
is the OpenIFS stable-BL turbulent mixing --- the ``long-tail'' stability functions / minimum stable diffusion whose
lineage runs Louis (1979) $\to$ Beljaars \& Viterbo (1998) --- \emph{orthogonal} to the vegetation fix.

\textbf{Two counters this diagnosis must own} (both align with the report's own ``compensating error'' worry).
\emph{(a) The mixing fix is a documented compensating error:} enhanced long-tail stable mixing is stronger than
LES/observations justify and is retained only because reducing it degrades the large-scale flow --- it improves 2\,m
temperature but deepens cyclones and harms low cloud (Sandu et al.\ 2013; Holtslag et al.\ 2013), so Phase~1B must
check the storm-track/low-cloud climatology (which CMPI/\texttt{part8} will not catch cleanly). \emph{(b) The inversion
may be a cloud-phase problem:} Pithan et al.\ (2014) attribute much of the modelled Arctic wintertime cold bias to
missing mixed-phase low cloud (a LW, cloud-\emph{phase} effect) --- our null control ruled out cloud \emph{fraction}
but not phase, so we scope the residual claim to ``given cloud fraction is unbiased'' and flag a cloud-phase /
LW$\downarrow$ test as a prerequisite. Likewise the $r=+0.76$ residual-to-inversion match is suggestive, not proof
that the ``fourth lever'' is truly orthogonal.

\begin{figure}[H]\centering
\includegraphics[width=0.99\textwidth]{plots/djf_second_driver.png}
\caption{Second-driver search: full-forest (no-LPJG) run DJFM vs ERA5, boreal NH. Residual cold (top-left) matches
the inversion-excess (top-right, $r=+0.76$); MSLP (bottom-left) and cloud (bottom-right) biases are diffuse and
uncorrelated.}
\end{figure}

\textbf{Near-complete attribution.} The two mechanisms close the boreal DJF cold bias between them: of the
$\approx-4.5$\,K Siberian bias vs ERA5, the LPJG forest-cover collapse (snow-albedo in spring, roughness in winter)
contributes $\approx-3$\,K, and the over-strong stable-BL inversion the remaining $\approx-1.4$\,K. Two
\emph{orthogonal} levers --- the LPJG establishment fix and enhanced stable-BL mixing --- therefore target
essentially the entire bias. (The veg cooling is also a \emph{stable} offset, not a runaway: over 20\,yr the LPJG
run carries more/thicker Arctic ice as a bounded mean offset with no amplifying sea-ice drift.)

\subsection{Working conclusion \& lever}
LPJ-GUESS sends essentially no boreal \emph{forest cover} to OpenIFS (cvh $\sim0.6\to0.02$ over Siberia,
\S\ref{sub:withwithout}), so the canopy no longer masks snow (spring albedo, \S\ref{sub:withwithout}) \emph{nor}
roughens the surface (winter inversion, \S\ref{sub:roughness}) --- a physically robust driver of the
Siberia/N.\ Canada cold bias, consistent with the EC-Earth family in which LPJ-GUESS is interactively coupled
(D\"oscher et al.\ 2022) and in which dynamic vegetation-cover / effective-LAI improves high-latitude spring climate
through the snow-albedo coupling (Alessandri et al.\ 2017). Source trace (verified): LPJG computes the high-veg fraction and
exports it over OASIS as \texttt{GUE\_FRAH}$\to$\texttt{cvh} (\texttt{framework.cpp::computeDailyLAIandFPCforIFS};
\texttt{namcouple} line 120). With the AWI setting \texttt{iftreefracca=1} (\texttt{global.ins}:151), cvh is the
summed tree \emph{crown area}$\times$density (no 0.05 clamp --- that constant is dead code; \texttt{natural\_frac}$=1$
over PI land), so cvh$\approx$0.05 means boreal tree cover is genuinely $\sim$5\%. The establishment gates
(\texttt{gdd5min\_est=500}, \texttt{twmin\_est=5}) are \emph{not} the blocker --- they are cleared across the boreal ---
so this is a \textbf{tree-vs-grass competition failure} (grass out-produces trees at canopy closure; settled in
\S\ref{sub:decisive}), with the degraded state \emph{inherited from the LPJG spin-up}, not a coupled runaway. Seasonally (\S\ref{sub:seasonality})
the LPJG cold effect spans September--June (DJF $-2.4$, MAM $-2.5$\,K over the boreal box)
and vanishes only in July--August --- i.e.\ it is a cold-season, not a growing-season, effect. The co-equal
\emph{sunless} late-autumn/winter part --- which pure snow-albedo cannot explain --- is now
resolved (\S\ref{sub:withwithout}--\ref{sub:residual}): it is the roughness$\to$inversion effect of the lost
forest (\S\ref{sub:roughness}) plus a model-intrinsic over-strong stable-BL inversion (\S\ref{sub:residual}).

\textbf{Fixes}: (1) tune the boreal \textbf{tree-vs-grass competition} (\texttt{greff\_min}, C3-grass competitiveness,
cold-tree photosynthesis \texttt{pstemp\_low}) --- \emph{not} the establishment gates, which are correct
(\S\ref{sub:decisive}) --- and re-do the LPJG spin-up; highest value, cheap to screen offline by checking cvh recovers
toward the MODIS pattern; (2) revisit \texttt{iftreefracca}/crown allometry; (3) a woody stem-area floor on the
snow-masking fraction in OpenIFS
(\texttt{surfbc\_ctl\_mod.F90}:389--391 / \texttt{surfrad\_ctl\_mod.F90}:650--655) --- but this only bites once (1)
raises cvh, since the masked-snow tile is weighted by cvh itself. Fixes (1)--(3) are the physically-grounded
\emph{warming lever} of \S\ref{sec:rec} and buy $\sim$3\,K (\S\ref{sub:levervsera5}). \textbf{A fourth, orthogonal
lever} attacks the residual $\sim$1.4\,K (\S\ref{sub:residual}): \emph{enhanced stable-boundary-layer mixing} in
OpenIFS (VDF stability functions / minimum stable diffusion), which would relax the over-strong surface inversion.
Together (1)--(4) target essentially the entire boreal DJF cold bias.

\subsection{Decisive test: competition, not climate \textnormal{---
\textcolor{red}{substantially qualified, see \S\ref{sub:xfer}}}}\label{sub:decisive}
\fbox{\parbox{0.97\textwidth}{\small\textbf{Read this section with \S\ref{sub:xfer} alongside.} The argument
below turns on a standalone LPJG re-spin driven by the OIFS--AMIP forcing under-vegetating despite the
establishment gates being cleared, which was read as an LPJG-internal competition failure. A later controlled
experiment (\S\ref{sub:xfer}) shows that \emph{moving between the CRUNCEP and AMIP forcings alone}, with
byte-identical parameters, costs $-54\,\%$ of Siberian and $-66\,\%$ of East-Siberian \texttt{TREEFPC}. The
under-vegetation seen here is therefore largely what the forcing change predicts on its own, and cannot by
itself demonstrate a competition failure. \dead{What does survive is FORESTFPC$=0$ --- no closed canopy forms
under either forcing --- so the closed-canopy failure remains unexplained and may still be competition.}
\texttt{FORESTFPC} is in fact a \emph{void diagnostic} here: it counts only the managed-forest land-use class,
which is empty in this PI configuration, so it is identically zero whatever the vegetation does
(\S\ref{sub:xfer}). There is no separate closed-canopy failure.
\dead{Verdict: the deficit is primarily an LPJG-internal competition problem, not a climate one.}}}
Is the boreal forest deficit a \emph{climate} problem (fix the atmosphere) or an \emph{LPJG-internal} tree-vs-grass
competition problem (fix the vegetation)? We answered it directly with a \textbf{feedback-free} experiment: a clean
OpenIFS--AMIP run (prescribed observed SST, no coupled ocean/LPJG cold feedback; $\sim$PI 1870--79) generated a
realistic daily forcing, which drove a 2000-yr standalone LPJG re-spin on a \emph{dense} 11\,538-cell gridlist
(build detail in \S\ref{sub:amipgen}). Fig.~\ref{fig:decisive} puts that forcing (top) beside the equilibrium
vegetation it produced (bottom). \textbf{The forcing is \emph{not} growing-season-cold}: over the Siberian box the
warmest-month T2m is $14.7$\,\textdegree C (100\% of cells clear the \texttt{twmin\_est}$=+5$ gate) and GDD5 is $730$
(76\% clear $500$) --- establishment is climatically permitted essentially everywhere trees should be. \textbf{Yet the
spin-up under-vegetates and grass wins}: cvh$=0.21$ (far below the $0.4$--$0.7$ closed-taiga range),
\texttt{GRASSFPC}$=0.37 > \texttt{TREEFPC}=0.21$, and \dead{no closed canopy forms anywhere (FORESTFPC$=0$)}\footnote{\textbf{Void diagnostic.} \texttt{FORESTFPC} sums only PFTs with \texttt{landcover==FOREST} (\texttt{commonoutput.cpp}:1143), i.e.\ the \emph{managed-forest} land-use class; natural trees are counted in \texttt{TREEFPC} via the \texttt{NATURAL} branch. In this PI configuration nothing is allocated to FOREST (\texttt{Natural\_sum}$=0.684$, \texttt{Total}$=0.665$), so \texttt{FORESTFPC}$\equiv0$ \emph{by construction} while \texttt{TREEFPC}$=0.323$. It is a land-use bookkeeping artifact and says nothing about canopy closure.}, with grass
filling exactly the East-Siberian interior. On a dense gridlist this cannot be dismissed as under-sampling.
\textbf{Verdict: the deficit is primarily an LPJG-internal competition problem, not a climate one} --- so the primary
counter-lever is the vegetation model (boreal \texttt{greff\_min}, C3-grass competitiveness, \texttt{pstemp\_low}), not
another atmospheric retune. The one climate co-limit is the deep-winter cold ($-28.3$\,\textdegree C box mean,
$<-40$ in NE Siberia) that gates out evergreens on $\sim40\%$ of cells --- but climatically-permitted \emph{larch}
(no cold gate) fails there too (BNS cover $0.035$), so the winter warming of Phase~1B is only a \emph{secondary} aid.

\begin{figure}[H]\centering
\includegraphics[width=0.99\textwidth]{plots/v3_decisive_climate_vs_veg.png}
\caption{\textbf{Decisive climate-vs-competition test} (NH boreal, equilibrium yr 3900). \emph{Top:} the
feedback-free obs-SST $\sim$PI OIFS--AMIP forcing that drove the standalone LPJG re-spin --- warmest-month T2m
(vs \texttt{twmin\_est}$=+5$), annual GDD5 (vs $500$), coldest-month T2m (vs evergreen \texttt{tcmin\_est}$=-30$).
\emph{Bottom:} the equilibrium vegetation it produced --- cvh, \texttt{TREEFPC}, \texttt{GRASSFPC}. The growing-season
forcing clears the gates across the boreal, yet the spin-up yields cvh$\approx0.2$ with grass dominating the
East-Siberian interior and no closed canopy --- grass out-competes trees under an adequate climate. The only climate
co-limit is the deep-winter cold (top-right, NE Siberia), but climatically-permitted larch fails there too.}
\label{fig:decisive}
\end{figure}

\part{The path forward: raise the floor, then tune}
\emph{Parts~I--II give the strategy: the imbalance can only be cut by cooling, and cooling is only affordable once the
structural NH cold floor is raised. So first fix the floor (vegetation competition + stable-BL mixing), \textbf{then}
spend a moderate cooling lever to reach surface balance --- now without the RMSE blow-up. First the caveats that
condition this, then the recommendation and a phased plan.}

\section{Caveats}\label{sec:caveats}
\begin{enumerate}\itemsep1pt
\item \textbf{06D ``HRlike'' has no HR-specific namelist change} --- only the common pond shift (identical to 06C);
its metrics confirm this. If a resolution/timestep/viscosity change was intended, it is not in the namelists.
\item \textbf{Filename convention.} These runs write \texttt{atm\_remapped\_1m\_<var>\_YYYY-YYYY.nc} (no
\texttt{\_1m\_} infix). I added backward-compatible fallbacks to \texttt{part2}, \texttt{part8} and the CMPI
preprocessing; existing configs are unaffected.
\item \textbf{Accumulation period} for these XIOS runs is \textbf{3600\,s} (not the 21600\,s of the old HR config);
configs use \texttt{accumulation\_period=3600} (verified: tsr\,$\to\approx240$\,W/m$^2$).
\item None of the runs is at equilibrium within 30\,yr --- numbers are a ranking, not final values.
\end{enumerate}

\section{Recommendation \& synthesis}\label{sec:rec}
Reading global CMPI/RMSD as a guide rather than a gate (the PI-vs-PD offset is expected), the picture is:
\begin{itemize}\itemsep1pt
\item The \textbf{melt-pond / albedo family (06A--06D, etc.)} are \emph{clean}: they improve sea ice where they
should and leave everything else essentially untouched. Their imbalance gain is modest ($5$--$20\%$) but it comes
with no unexpected collateral --- safe building blocks.
\item The \textbf{ocean-mixing / coupling family (06O/T)} delivers the only large imbalance reductions
($30$--$38\%$) \emph{and} genuinely improves the ocean interior and sea ice --- but the cooling leaks into the
atmosphere globally (incl.\ tropics), which is degradation in places the parameter should not reach. That makes
them promising but not yet ``free.''
\item \textbf{06V} pushes too far: the \texttt{ENTSTPC3} convection change degrades \texttt{tas}/\texttt{zg}
everywhere --- not recommended as-is.
\end{itemize}
Best-balanced single steps: \textbf{06P} ($-29\%$, well-localised) and \textbf{06H} ($-20\%$).

\textbf{Synthesis --- and where the warming lever comes from.} The two arcs join here. The imbalance can only be
cut substantially by cooling levers (ocean mixing/coupling), so every combo tried so far overshoots into cold
bias. But \S\ref{sec:lpjg} shows the shared Siberia/N.\ Canada cold lobe is \emph{not} a tuning artefact, and a
purpose-built clean no-LPJG CORE3 reference (\S\ref{sub:withwithout}) now \emph{quantifies} it: of the $\approx-4.5$\,K
Siberian DJF bias vs ERA5, $\approx-3$\,K is the LPJG forest-cover collapse (via snow-albedo in spring \emph{and}
surface roughness in winter, \S\ref{sub:roughness}) and $\approx-1.4$\,K is a model-intrinsic over-strong stable-BL
inversion (\S\ref{sub:residual}). The clean run is \emph{better} than the Baseline globally (RMSD 1.514 vs 1.717\,K)
at \emph{near-unchanged} imbalance ($+0.2$\,W/m$^2$, within drift) --- indicating both are \emph{warming levers
largely orthogonal to} the TOA tuning. That reframes
the path forward:
\begin{enumerate}\itemsep1pt
\item Fix the LPJG boreal forest-cover bug (establishment limits \texttt{gdd5min\_est}/\texttt{twmin\_est} and/or
zero-LAI stem snow-masking --- see the companion note) and re-do the spin-up; this warms the NH high latitudes in
\emph{every} run, lowering the shared RMSD floor by $\approx$3\,K in the worst region.
\item Attack the residual with \emph{enhanced stable-BL mixing} in OpenIFS (\S\ref{sub:residual}) --- an
independent knob that relaxes the surface-trapped winter inversion.
\item On that warmer baseline, apply a \emph{moderate} mixing/coupling setting (06O-like, not 06V) for the
imbalance, plus the clean melt-pond changes --- the cooling it adds is then offset rather than additive.
\item Re-judge on the \emph{pattern} of CMPI change (\S\ref{sub:withwithout} method), not the global value, since
the PI-vs-PD offset is expected.
\end{enumerate}
In short: don't chase the imbalance with ever-colder atmospheres; remove the structural cold bias first, then the
moderate cooling needed for the imbalance becomes affordable.

\subsection{2026--07--24 coupled follow-up from the CRUNCEP3-initialized branch: the year-50 land--ocean split refines the path}\label{sub:cruncep3_init_followup}
After the earlier report structure was drafted, the coupled
CRUNCEP3-initialized branch (alias ``080a'') reached
50 years, which adds an important constraint on the tuning strategy. The new
result is \emph{not} a simple hemisphere-wide cold drift. By year 50 the system
shows a marked \textbf{warm-ocean / cold-land split}, while boreal tree cover
continues to retreat --- most strongly over Siberia / East Siberia.

Two points matter:
\begin{enumerate}\itemsep1pt
\item The older LPJG story remains necessary: the spin-up / competition state is
still part of the problem, and the vegetation levers are real.
\item But the 50-year coupled adjustment shows that a \emph{regional continental-land}
climate failure also develops, so vegetation-only retuning is not sufficient.
\end{enumerate}

The finished year-50 diagnostics make the asymmetry explicit. Relative to year 0,
NH45+ \texttt{TREEFPC} declines from $0.336$ to $0.215$, but the strongest loss is
concentrated over Siberia / East Siberia. In the Siberian box (55--75\textdegree N,
60--180\textdegree E) \texttt{TREEFPC} drops by $-0.188$, and in East Siberia
(55--75\textdegree N, 90--160\textdegree E) by $-0.232$, accompanied by a strong
drop in \texttt{AGDD5} ($-89$ and $-133$ respectively). At the same time,
NH45+ aggregate \texttt{AGDD5} actually \emph{increases}, and Scandinavia / Canada
also warm in thermal-growing metrics. So the problem is not ``the NH is too cold'';
it is that \textbf{boreal continental land, especially East Siberia, remains too
cold relative to the oceanic state}. This is exactly what one sees in the year-50
seasonal-temperature maps against the CRUNCEP3 reference: ocean-adjacent sectors
are warm-biased while large boreal land sectors remain cold.

\begin{figure}[H]
\centering
\includegraphics[width=0.98\textwidth]{plots/treefpc_coupled_cruncep3_init_y50_minus_y0_nh.png}
\caption{Coupled CRUNCEP3-initialized branch year-50 follow-up: NH \texttt{TREEFPC} at year 0 (1350),
final year (1399) and change (Y50$-$Y0). The retreat is broad, but the largest
loss is concentrated over Siberia / East Siberia.}
\end{figure}

\begin{figure}[H]
\centering
\begin{subfigure}{0.49\textwidth}
\includegraphics[width=\textwidth]{plots/t2m_seasonal_bias_coupled_cruncep3_init_vs_cruncep3_nh.png}
\caption{Early coupled decade: 1350--1359 vs CRUNCEP3 forcing}
\end{subfigure}
\hfill
\begin{subfigure}{0.49\textwidth}
\includegraphics[width=\textwidth]{plots/t2m_seasonal_bias_coupled_cruncep3_init_1390_1399_vs_cruncep3_nh.png}
\caption{Final coupled decade: 1390--1399 vs CRUNCEP3 forcing}
\end{subfigure}
\caption{Seasonal NH T2m bias of the coupled CRUNCEP3-initialized branch against CRUNCEP3 forcing. The
year-50 state is not a uniform NH cooling; instead it shows a land--ocean split,
with boreal continental land still cold while oceanic sectors are relatively warm.}
\end{figure}

This follow-up changes the recommendation in one practical way: the next runs
should \textbf{fix the physical coupled branch first, and only then add LPJ
compensation}. In other words, the "warm-floor" concept of this report still
holds, but the relevant floor is now more specific: \textbf{a warmer boreal land
floor}, not merely a warmer NH mean. The cleanest next matrix is therefore:
\begin{enumerate}\itemsep1pt
\item \textbf{06T-only} physical branch:
	\texttt{oasis3mct.time\_step=3600}, \texttt{use\_momix=True},
	\texttt{momix\_kv=0.01}, \texttt{kpp\_av0=0.003}, \texttt{kpp\_kv0=0.003},
	\texttt{kpp\_kvbckg=5e-06}.
\item \textbf{06V-only} physical branch: 06T plus \texttt{ENTSTPC3=1}.
\item Best of 06T/06V plus \textbf{07A} (\texttt{BNS greff\_min=0.035}).
\item Best of 06T/06V plus \textbf{07A+07C}
	(\texttt{C3G pstemp\_low=12}).
\end{enumerate}

The year-50 evidence also explains why \textbf{LPJ-only tuning should not be the
first move}. The all-tuned coupled NH bias panels (seasonal and annual) show that
the coupled atmosphere response already differs strongly among branches even before
the LPJ compensation is evaluated. If the physical branch leaves East Siberia cold,
then LPJ-only changes risk compensating the wrong climate state rather than fixing
it. Conversely, if a 06T/06V-type branch can hold the global net-radiation target
while reducing the East-Siberian land cold pool, then the LPJ levers can be used
for their intended role --- retaining boreal trees --- rather than masking a
regional coupled-climate defect.

\begin{figure}[H]
\centering
\includegraphics[width=0.98\textwidth]{plots/net_radiation_panel_baseline_A_to_V_07A_to_C_HRgoal_LRpi_last30.png}
\caption{Global net-radiation history used to define the next physical branch.
The preferred candidates remain the late 06-series branches that move closest to
the target while staying within a plausible tuning envelope; the follow-up now
adds the requirement that they must also improve boreal land climate, not only the
global radiation budget.}
\end{figure}

\begin{figure}[H]
\centering
\begin{subfigure}{0.49\textwidth}
\includegraphics[width=\textwidth]{plots/t2m_bias_vs_base06_all_tuned_coupled_nh_10yrseasonal_DJF.png}
\caption{DJF}
\end{subfigure}
\hfill
\begin{subfigure}{0.49\textwidth}
\includegraphics[width=\textwidth]{plots/t2m_bias_vs_base06_all_tuned_coupled_nh_10yrseasonal_MAM.png}
\caption{MAM}
\end{subfigure}
\par\medskip
\begin{subfigure}{0.49\textwidth}
\includegraphics[width=\textwidth]{plots/t2m_bias_vs_base06_all_tuned_coupled_nh_10yrseasonal_JJA.png}
\caption{JJA}
\end{subfigure}
\hfill
\begin{subfigure}{0.49\textwidth}
\includegraphics[width=\textwidth]{plots/t2m_bias_vs_base06_all_tuned_coupled_nh_10yrseasonal_SON.png}
\caption{SON}
\end{subfigure}
\caption{All tuned coupled experiments relative to 06\_Baseline, resolved by
season (10-year means). These panels help separate branches that reduce the global
radiation problem by globally cooling everything from branches that have a more
usable boreal-land signal.}
\end{figure}

\section{Work plan: a systematic path to a warm-floor, surface-balanced PI control}\label{sec:workplan}
\textbf{Two revised targets.} (i) net \emph{surface} energy balance $|F_\mathrm{sfc}|\lesssim0.3$\,W/m$^2$ ---
\emph{not} TOA. Surface and TOA net flux track to within $\approx0.1$\,W/m$^2$ in every run (\S\ref{sec:rad}); they
are one problem, so a surface-zero target is equivalent to TOA\,$\approx-0.1$\,W/m$^2$, marginally stricter than
TOA$=0$. (ii) global 2\,m-T bias vs ERA5 $\approx-1.3$\,\textdegree C --- the physically-expected PI-minus-present-day
offset, \emph{not} zero: a run that reaches $-1.3$\,\textdegree C is \emph{correct}, not cold-biased.

\textbf{Ordering principle.} The imbalance can only be cut by cooling, and every cooling lever over-cools if applied
to today's cold-biased base. So we \emph{first} raise the warm floor with two structural fixes that are
\emph{orthogonal to the TOA balance} (they change the land surface and BL mixing, not the radiation budget ---
\S\ref{sub:withwithout} showed the no-LPJG floor sits at \emph{near-unchanged} imbalance, $+0.2$\,W/m$^2$ within
drift), \emph{then} spend a
\emph{moderate} cooling lever to reach surface balance, now affordable because the floor is $\sim3$--$4$\,K warmer
where it was worst. Phases 1A/1B are independent (parallel); Phase 2 depends on both; Phase 3 gates.

\subsection{Phase 1A --- Restore boreal tree cover (LPJG establishment)}
\textbf{Literature \& observational basis --- what realistic trees require (so we don't green the tundra).} The
observed arctic/boreal treeline sits near the July $10$\,\textdegree C isotherm (K\"oppen), refined by K\"orner \&
Paulsen (2004) to a growing-season mean of $\sim6$--$7$\,\textdegree C over $\sim90$ days, with the northern boreal
limit at $\lesssim30$ days $>10$\,\textdegree C. \textbf{Crucially, LPJG's current gates are already the standard,
treeline-calibrated values} (Smith et al.\ 2001; Sitch et al.\ 2003 --- byte-identical to the canonical
\texttt{global.ins}): BNE/BINE \texttt{twmin\_est}$=5$ (already \emph{below} the $10$\,\textdegree C isotherm) and
\texttt{gdd5min\_est}$=500$ sit essentially \emph{at} the observed treeline --- so loosening them risks trees north
of the real treeline. Two further realism facts constrain the fix: (i) NE Siberia's forest is $\sim$80\,\% \emph{larch}
(deciduous \texttt{BNS}), because evergreen conifers cannot survive its extreme winters (Kharuk et al.; Herzschuh
2020, \emph{Glob.\ Ecol.\ Biogeogr.}\ 29, 198--206, the deciduous-vs-evergreen distribution paper) --- LPJG already
encodes this (BNE/BINE \texttt{tcmin\_surv}$=-31$ excludes them; BNS has \emph{no} cold gate and
\texttt{gdd5min\_est}$=350$), so the realistic tree there is larch, and it is \emph{already permitted}. (ii) Closed
taiga tree cover is $\sim0.4$--$0.7$ (MODIS VCF), falling below $0.2$ at the taiga--tundra ecotone --- the spatial
shape cvh must reproduce.

\textbf{Larch caveat.} The realistic NE-Siberian tree is \emph{deciduous} larch, whose leafless winter canopy masks
little snow --- so the spring \emph{albedo} mechanism is genuinely weaker over the larch zone (though it retains full
force over the evergreen west Eurasia / N.\ America). This only re-weights the two mechanisms, it does not soften the
conclusion: the \emph{roughness} mechanism survives (even leafless, larch is $\gg$ rougher than grass/tundra), and the
seasonality fingerprint (\S\ref{sub:seasonality}) shows the LPJG cold effect is co-equal in the \emph{sunless}
Nov--Jan months --- i.e.\ the dark-season roughness/inversion part is what carries the larch zone.

\textbf{The coupled winter cold is a self-inflicted trap, not the trigger.} At the cells where prescribed veg has
forest but LPJG does not (1785 boreal-land cells), the growing-season gates are already cleared in every run
(JJA $\ge12$\,\textdegree C, GDD5 $>500$); only the winter \texttt{tcmin\_est}$=-30$ evergreen gate binds. But against
the \emph{forested} WITHOUT-LPJG climate it excludes evergreens on just $23\%$ of cells ($\approx$ERA5's $22\%$),
vs $43\%$ in the degraded WITH-LPJG run --- the missing forest's own $-3.3$\,\textdegree C winter cooling doubles the
exclusion. So it is a \textbf{bistable trap} seeded by the under-forested spin-up, not a climate trigger.

\textbf{The decisive test settles it (\S\ref{sub:decisive}): competition, not climate} --- so Phase 1A's lever is the
\textbf{competition parameters}, not the establishment gates: boreal \texttt{greff\_min} (0.03 BNE / 0.09 BINE,BNS),
C3-grass competitiveness, and cold-tree photosynthesis \texttt{pstemp\_low}$=10$\,\textdegree C. Robustly, trees
\emph{do} establish (nonzero density) and fire is never involved (burnt fraction $=0$ despite BLAZE) --- the failure
is canopy-closure/productivity competition, not a gate or a fire problem. Phase-1B winter warming
(\S\ref{sub:phase1b}) is a \emph{secondary} aid, only on the coldest evergreen-gated cells.

\textbf{Why not just relax the gates?} The summer gates are not binding (above), so loosening them cannot help; and
the \emph{one} binding gate --- the winter \texttt{tcmin\_est} --- correctly excludes evergreens from the
extreme-cold larch zone, so relaxing it would plant \emph{unrealistic} evergreens where even ERA5 says larch. The
fix is therefore the winter \emph{climate}, not the limits. (Keep \texttt{iftreefracca 1}, \texttt{global.ins}:150:
cvh is the summed tree \emph{crown area}, so it tracks establishment success directly.)

\textbf{Screen \& deploy.} Tune the competition knobs on the cheap standalone LPJG spin-up (the OIFS--AMIP forcing of
\S\ref{sub:decisive}; minutes--hours of wall-time), targeting the MODIS-VCF pattern ($\sim0.4$--$0.7$ closed taiga,
$<0.2$ at the ecotone, larch winning the cold NE) and \dead{a closed canopy (FORESTFPC$>0$)}\footnote{\textbf{Void diagnostic.} \texttt{FORESTFPC} sums only PFTs with \texttt{landcover==FOREST} (\texttt{commonoutput.cpp}:1143), i.e.\ the \emph{managed-forest} land-use class; natural trees are counted in \texttt{TREEFPC} via the \texttt{NATURAL} branch. In this PI configuration nothing is allocated to FOREST (\texttt{Natural\_sum}$=0.684$, \texttt{Total}$=0.665$), so \texttt{FORESTFPC}$\equiv0$ \emph{by construction} while \texttt{TREEFPC}$=0.323$. It is a land-use bookkeeping artifact and says nothing about canopy closure.} --- tree NPP $>$ grass NPP ---
\emph{not} merely a higher global mean. On a passing screen, re-do the coupled LPJG spin-up, regenerate the restart,
and carry it into Phase 2.

\subsection{Execution: the clean AMIP--noLPJG forcing generator}\label{sub:amipgen}
The decisive test (\S\ref{sub:decisive}) needs a \emph{feedback-free} climate to drive it --- reusing coupled-LPJG
output would carry the very cold bias we are testing. We therefore built a dedicated \textbf{OpenIFS--AMIP} generator
at TCO95 (\texttt{oifsamip-cy48}, git-aligned to \texttt{awiesm3-develop}): AMIP prescribes observed SST/sea-ice, so
the atmosphere has no coupled ocean/LPJG feedback. Building and running it required fixing three esm-tools bugs (a
FESOM-less crash, missing XIOS/coupling compile flags, a missing WAM date-preprocessing recipe --- all committed \&
pushed; detail in App.~\ref{app:impl}). A 10-yr $\sim$PI decade (1870--79 observed SST, GHG fixed 1850) produced the
daily forcing. \emph{One methodological catch worth recording:} the first re-spin came out barren everywhere because
the XIOS field-def de-accumulated fluxes over the 6\,h \texttt{freq\_op} cadence instead of the true
\texttt{NFRHIS$\times$TSTEP}$=1$\,h period, making \texttt{rsns}/\texttt{rlns}/\texttt{pr} exactly 6$\times$ too low
(net SW $28$ vs $166$\,W/m$^2$) --- fixed by a validated $\times6$ patch and, at root, by templating the divisor from
esm-tools. \textbf{Expected payoff} of the whole Phase~1A (\S\ref{sub:levervsera5}): $\approx-3$\,K off the Siberian
DJF bias, $\approx-0.2$\,K global RMSD, at near-unchanged imbalance.

\subsection{Phase 1B --- Enhance stable-boundary-layer mixing (OpenIFS)}\label{sub:phase1b}
\textbf{Change.} The residual $\approx1.4$\,K surface-trapped inversion (\S\ref{sub:residual}) is relaxed by
strengthening stable-BL turbulent mixing in the EDMF K-diffusion scheme via the ``long-tail'' (LTG) stability
functions --- lengthening the tail keeps mixing non-zero at large Richardson number and weakens the over-strong
surface inversion. There is no namelist for the stable branch, so this is $\sim$2 single-line OpenIFS source edits
(minimum stable mixing length \texttt{ZLMIN} and the LTG coefficient \texttt{ZCB}) plus recompile; the exact
file:line targets and secondary companions are in App.~\ref{app:impl}. \emph{Bear the documented side-effects in
mind} (\S\ref{sub:residual}): long-tail mixing beyond the observational range can deepen cyclones and degrade low
cloud, so verify the storm-track/low-cloud climatology, and pre-test cloud phase / LW$\downarrow$ (Pithan et al.\
2014) since the residual may be partly a cloud-microphysics rather than turbulence error.

\textbf{Test.} Change \emph{one} knob first (\texttt{ZLMIN}) to keep attribution clean; scan 2--3 values with short
(1--5\,yr) coupled sensitivity runs of the no-LPJG reference. \textbf{Gate:} the DJFM low-level inversion bias
(\texttt{djf\_second\_driver.py}) moves from $+1.4$\,K toward ERA5's $+4.2$\,K level and the residual boreal cold
shrinks, \emph{without} eroding legitimate inversions or warming the tropics (judge globally via \texttt{part8}/CMPI
pattern). \textbf{Risk:} over-mixing warms too broadly / degrades low cloud --- mitigate with small steps and
global CMPI-pattern judging.

\subsection{Phase 2 --- Moderate imbalance tuning on the warm floor}
With the floor raised (1A$+$1B) and global skill already better than Baseline (RMSD $1.514$, CMPI $0.738$;
\S\ref{sub:withwithout}), spend a \emph{moderate} cooling lever to bring the \emph{surface} budget to zero. Build the
combined config in layers, cleanest-first:
\begin{center}\small
\begin{tabular}{l l r r l}
\toprule
layer & lever(s) & $\Delta$TOA [W/m$^2$] & raw bias [K] & rationale \\
\midrule
base      & melt-pond (06A--06C; opt.\ 06G) & $-0.08$ to $-0.26$ & $-1.0$ to $-1.3$ & collateral-free \\
localised & 06P (\texttt{ENTSTPC3})         & $-0.44$ ($-29\%$)  & $-1.36$ & best efficiency $\sim1.05$\,W/m$^2$/K \\
moderate  & 06O (\texttt{1hcpl}+\texttt{mospp}) & $-0.55$ ($-36\%$) & $-1.77$ & improves ocean/ice; \emph{not} 06V \\
\midrule
(spare)   & $+$\texttt{kpplow} $=$ 06T      & $-0.58$ ($-38\%$)  & $-1.79$ & if more cut needed \\
\bottomrule
\end{tabular}
\end{center}
\textbf{Why this stops short of 06V.} 06V (\texttt{entstpc3}$+$\texttt{kpplow} all stacked) gives the biggest cut
($-57\%$) but overshoots to $-2.1$\,K and degrades \texttt{tas}/\texttt{zg}/\texttt{rlut} in \emph{every} region
including the tropics. The base$+$06P$+$06O stack cuts $\sim1.0$--$1.1$\,W/m$^2$ of surface imbalance (reference
surface $+1.52\to\sim0.4$--$0.5$; iterate the coupling strength or add \texttt{kpplow} to close the last $\sim0.4$).
\textbf{Cold-bias accounting:} the $\sim-1.7$\,K raw cooling of the coupling lever, applied on a floor lifted
$\sim3$--$4$\,K where the bias was worst (and $\sim0.2$--$0.4$\,K globally), lands the global mean near the target
$-1.3$\,\textdegree C rather than the $-1.8$ to $-2.1$\,\textdegree C the same lever produces today.

\subsection{Phase 3 --- Verification protocol and decision gates}
Run the combined config 30\,yr (evaluate the last decade). Judge with the \texttt{reval} toolchain \emph{plus} the DJF
driver diagnostics:
\begin{center}\small
\begin{tabular}{l l l}
\toprule
target & metric (tool) & accept \\
\midrule
surface balance   & \texttt{part2} ``Surface total''             & $|F_\mathrm{sfc}|\lesssim0.3$\,W/m$^2$ (TOA\,$\approx-0.1$) \\
cold bias         & \texttt{part8} 2\,m bias vs ERA5            & $-1.3\pm0.3$\,\textdegree C; RMSD $\le$ ref.\ $1.51$ \\
multivariate skill& \texttt{part4} CMPI                          & $\le$ Baseline $0.745$; \emph{pattern}, not just global \\
no new regionals  & CMPI by region (esp.\ tropics)             & no family-typical global cooling leak \\
driver check      & \texttt{djf\_second\_driver.py} / \texttt{...roughness...} & inversion bias $\to$ ERA5; cvh restored \\
\bottomrule
\end{tabular}
\end{center}
\textbf{Iteration \& critical path.} Each axis adjusts independently (the floor fixes are orthogonal to the imbalance
lever): too cold $\to$ back off coupling; imbalance too high $\to$ add \texttt{kpplow}; inversion too strong $\to$ raise
\texttt{ZLMIN}. Phases 1A and 1B run in parallel now (both cheap/offline); 1A's coupled re-spin is the long pole, and
the Phase-2 run starts once the re-spun restart and the chosen stable-BL build are both in hand.

\vspace{4pt}
{\footnotesize All scripts, data, figures and this report live in \texttt{/work/bb1469/a270092/eval/}\,:
\texttt{data/MASTER\_metrics.csv}, \texttt{data/cmpi\_by\_variable.csv}, \texttt{data/cmpi\_by\_region.csv}, \texttt{configs/},
\texttt{plots/}, \texttt{run\_figures/<run>/}, \texttt{notebooks/}. Tool: \texttt{release\_evaluation\_tool2}
(reval env); per-run diags via \texttt{part2\_rad\_balance.py} / \texttt{part8\_t2m\_vs\_era5.py} /
\texttt{part4\_cmpi.py}.}

\clearpage
\part{The atmosphere alone --- attribution, and where the remaining error lives}
\emph{Parts~I--III argued from coupled runs, in which atmosphere, ocean and vegetation errors are
inseparable. This part removes that ambiguity with a \textbf{feedback-free OpenIFS--AMIP reference}:
prescribed observed SST/sea-ice, prescribed vegetation, same OIFS branch as the coupled campaign.
Two results follow. \textbf{(i)} With SST prescribed, the model still runs $+0.67$\,W/m$^2$ out of TOA
balance --- that fraction of the imbalance cannot be ocean heat uptake and is a genuine atmospheric
radiative error, which places a \emph{floor} beneath what any ocean-side lever can achieve.
\textbf{(ii)} The boreal \emph{summer} cold bias survives in full without ocean or LPJ-GUESS, so it is
the atmosphere's --- and it is a surface \emph{shortwave} deficit, not the boundary-layer problem of
\S\ref{sub:residual}. Together these redirect the work plan: the next tuning should happen on AMIP,
where both remaining errors live and where an iteration costs $\sim$7$\times$ less than a coupled run.}

\section{The feedback-free AMIP reference (2026--07--27/28)}\label{sec:amip}
\texttt{amip\_pi\_base}, \texttt{/work/bb1469/a270092/runtime/oifsamip-cy48/}: OpenIFS TCO95L91, GHG/ozone/
aerosol fixed at 1850 (\texttt{NCMIPFIXYR=1850}), driven by the earliest observed AMIP SST decade
1870--1879 --- a self-consistent $\sim$PI atmosphere directly comparable to the PI-control coupled runs.
10\,yr complete in 2:44 on 14 nodes.

Two configuration decisions make it a valid counterpart to the coupled model, and both were wrong in the
earlier \texttt{\_lpjgforcing} run that Part~II used (\S\ref{sub:amipgen}):
\begin{itemize}\itemsep1pt
\item \textbf{OIFS branch.} The oifsamip tree was on \texttt{local\_combined\_fixes}\,@\,\texttt{1303a90};
the coupled campaign runs esm-tools version \texttt{48r1v5} $=$ \texttt{movcav-landice+co2-concdriven}.
The branches differ in \emph{atmospheric physics}, not merely coupling plumbing (\texttt{suecozv},
\texttt{updecozv}, \texttt{suecrad}, \texttt{suphec}, \texttt{srfi\_mod}, \texttt{surfece},
\texttt{voskin\_mod}). Rebuilt at \texttt{f3ccacb} --- the last non-debug commit, and the same state as
our coupled tree.
\item \textbf{Sea albedo.} The old runscript carried \texttt{TMPRALBSEAD}$=0.075$. The coupled runs never
set it, so they use the source default \texttt{RALBSEAD}$=0.06$ (\texttt{susrad\_mod.F90}:107); OASIS
exchanges only the \emph{ice} albedo (\texttt{sia\_feom}$\to$\texttt{A\_Ice\_albedo}), so OpenIFS computes
open-water albedo itself in coupled mode too. A $0.015$ albedo offset over ocean is worth order
$1$\,W/m$^2$ of reflected SW --- comparable to the quantity being attributed. Removed.
\end{itemize}

\subsection{Which budget the target applies to: snow enthalpy}\label{sub:snowenth}
\texttt{part2\_rad\_balance.py}:255--268 computes
\[
F_\mathrm{sfc}=\mathrm{ssr}+\mathrm{str}+\mathrm{sshf}+\mathrm{slhf}-\mathrm{sf}\cdot\rho L_f,
\qquad F_\mathrm{TOA}=\mathrm{tsr}+\mathrm{ttr}.
\]
The snowfall-enthalpy term enters the \emph{surface} budget only, and correctly so: it is an internal
atmosphere$\leftrightarrow$surface transfer and cancels at the top of the atmosphere. Every net-TOA number
quoted in this report is $(\mathrm{tsr}+\mathrm{ttr})/\Delta t_\mathrm{acc}$, i.e.\ exactly part2's TOA total.

\begin{center}\small
\begin{tabular}{l r r r r r}
\toprule
 & TOA & SFC w/o snow & snow enthalpy & SFC (part2) & TOA$-$SFC \\
\midrule
AMIP 1872--79 & 0.667 & 1.441 & 0.925 & \textbf{0.516} & $+0.152$ \\
09C 1370--79  & 0.918 & 1.897 & 0.944 & \textbf{0.953} & $-0.035$ \\
080a 1370--79 & 1.682 & 2.588 & 0.823 & \textbf{1.765} & $-0.083$ \\
\bottomrule
\end{tabular}
\end{center}

The term is $0.82$--$0.94$\,W/m$^2$ --- \emph{larger than the entire remaining gap to target}. Without it the
surface budget misses TOA by $\approx1$\,W/m$^2$; with it, 080a closes at $-0.083$, reproducing the
$\approx-0.09$ quoted in \S\ref{sec:rad}. \textbf{Convention warning:} the work plan's acceptance target is
on the \emph{surface} (\S\ref{sec:workplan}), where the same runs read $+0.52$ (AMIP), $+0.95$ (09C),
$+1.77$ (080a). State which convention is in use whenever these numbers are compared.
\emph{Open item:} AMIP's closure ($+0.152$) is 2--4$\times$ worse than the coupled runs'. Small against the
$+0.67$ signal, so the attribution is unaffected, but unexplained --- suspect the prescribed SST/sea-ice
surface tile. Chase it before quoting the AMIP floor to better than $\sim0.15$\,W/m$^2$.

\subsection{Result I --- roughly 45\,\% of the imbalance is atmospheric}\label{sub:amiptoa}
\textbf{AMIP net TOA $=+0.67$\,W/m$^2$} (1872--79, after discarding two soil spin-up years; per-year
$0.13$--$0.92$; \texttt{tsr}$=240.3$\,W/m$^2$ confirms the 3600\,s accumulation divisor). With SST prescribed
this \emph{cannot} be ocean heat uptake.

Against a coupled baseline of $+1.44$ (080a), that is $\approx45\,\%$ of the coupled imbalance. The best
branch in the entire campaign (09C, $+0.92$) now sits only $0.25$\,W/m$^2$ above this floor: of the
$1.08$\,W/m$^2$ still separating it from the HR goal, $\approx0.83$ is atmospheric and only $\approx0.25$ is
reachable by ocean-side tuning. \textbf{This is the campaign's most consequential number} --- it says the
lever family of Part~I is nearly exhausted.

\emph{Caveat.} AMIP runs over \emph{observed} SST/ice while the coupled model runs over its own biased
state, so $+0.67$ is ``what this atmosphere does over a correct ocean'', not a strict lower bound on the
coupled value. But for a $\sim$PI configuration, where the real world was near equilibrium, the expected
AMIP TOA is $\approx0$ --- so most of the $+0.67$ is model error.

\subsubsection{Where does it live? Nowhere in particular --- the global budget is already right}\label{sub:toadecomp}
Decomposing the $+0.67$ against CERES-EBAF (\texttt{amip\_toa\_decomp.py}) was expected to localise it. It does
the opposite:

\begin{center}\small
\begin{tabular}{l r r r @{\hskip 1.2em} l r r}
\toprule
AMIP global & model & CERES & diff & band & net TOA & contribution \\
\midrule
absorbed SW      & 240.13 & 240.5 & $-0.37$  & 90S--60S        & $-90.5$ & $-6.06$ \\
OLR              & 239.46 & 240.2 & $-0.74$  & 60S--30S        & $-22.1$ & $-4.04$ \\
planetary albedo & 0.2947 & 0.293 & $+0.0017$& tropics 30S--30N& $+42.6$ & $+21.30$ \\
SW CRE           & $-44.40$ & $-45.4$ & $+1.00$ & 30N--60N     & $-21.8$ & $-3.98$ \\
LW CRE           & $+24.88$ & $+25.8$ & $-0.92$ & 60N--90N     & $-97.7$ & $-6.55$ \\
\textbf{net CRE} & $\mathbf{-19.52}$ & $\mathbf{-19.6}$ & $\mathbf{+0.08}$ & \textbf{sum} & & $\mathbf{+0.67}$ \\
\bottomrule
\end{tabular}
\end{center}

The cloud radiative effect matches observations to \textbf{$0.08$\,W/m$^2$} and the planetary albedo to
$0.002$. The AMIP and 09C zonal profiles lie almost on top of one another: the coupled model's radiation
\emph{is} the atmosphere's.

\dead{The band contributions reproduce the textbook meridional structure with no band carrying an anomalous
imbalance, so the Hyder et al.\ atmospheric-SW hypothesis does not appear as a gross cloud error and there is
no region to fix.} \textbf{That conclusion was an artifact of the comparison.} The bands above were compared
to \emph{each other}, and only the \emph{global} CRE to CERES. Differencing each band \emph{against CERES}
(\S\ref{sub:sobands}) shows the global agreement is a \textbf{cancellation of large regional errors}, and that
the Southern Ocean carries most of the imbalance.

\begin{figure}[H]\centering
\includegraphics[width=0.99\textwidth]{plots/amip_toa_decomposition_zonal.png}
\caption{Zonal-mean net TOA and cloud radiative effect, AMIP (atmosphere only) vs coupled 09C. The curves are
nearly indistinguishable, and the imbalance is the small residual of a large tropical surplus against
extratropical deficits --- not a localised error.}
\end{figure}

\subsubsection{Differenced against CERES per band: the Southern Ocean carries it}\label{sub:sobands}
\begin{center}\small
\begin{tabular}{l r r r r r}
\toprule
band & model & CERES & diff & area & \textbf{contribution} \\
\midrule
90S--65S              & $-97.0$ & $-104.0$ & $+7.03$ & $4.8\,\%$ & $\mathbf{+0.34}$ \\
\textbf{SO 65S--45S}  & $-49.5$ & $-55.3$  & $\mathbf{+5.81}$ & $9.8\,\%$ & $\mathbf{+0.57}$ \\
45S--30S              & $-5.6$  & $-1.4$   & $-4.20$ & $10.4\,\%$ & $-0.44$ \\
tropics 30S--30N      & $+42.6$ & $+45.1$  & $-2.51$ & $50.0\,\%$ & $-1.26$ \\
30N--45N              & $-4.8$  & $-5.7$   & $+0.89$ & $10.4\,\%$ & $+0.09$ \\
45N--65N              & $-50.1$ & $-54.1$  & $+4.02$ & $9.8\,\%$  & $+0.40$ \\
65N--90N              & $-106.3$& $-106.1$ & $-0.15$ & $4.8\,\%$  & $-0.01$ \\
\bottomrule
\end{tabular}
\end{center}

\textbf{The global budget is right by cancellation, not by being right.} The Southern Ocean
(65--45\textdegree S) runs $+5.8$\,W/m$^2$ too much absorbed energy and contributes
$\mathbf{+0.57}$\,W/m$^2$ to the global mean --- on its own, comparable to the entire $+0.67$ imbalance. The
cause is a cloud \emph{deficit}: cloud area $83.1\,\%$ against CERES's $89.7\,\%$ ($-6.6$\,pp) with SW CRE
$-60.3$ against $-68.1$ ($+7.8$\,W/m$^2$ too little reflection). \textbf{This is precisely the Hyder et al.\
(2018) signature}, and it vindicates the compensating-error worry of \S\ref{sec:so}: the ocean-mixing levers
of Part~I have been removing a Southern-Ocean warm bias that is substantially \emph{atmospheric} in origin.

\textbf{Consequence for the plan.} The $+0.67$ is not irreducible after all --- but neither is it fixable by a
global knob, since the global mean is already correct. It is fixable \emph{regionally}, and the two largest
regional errors have \emph{opposite} sign at TOA, which is what makes the round-10 design work
(\S\ref{sec:round10}). Re-framing the campaign arithmetic:
\begin{center}\small
\begin{tabular}{l r c r}
\toprule
 & net TOA & $=$ atmosphere $+$ & coupled-state \\
\midrule
080a (untuned coupled) & 1.44 & $0.67\ +$ & $\mathbf{0.77}$ \\
09C (06V $+$ newSeaIce) & 0.92 & $0.67\ +$ & $\mathbf{0.25}$ \\
\bottomrule
\end{tabular}
\end{center}
The ocean-side tuning has already delivered about two-thirds of what was available to it. Whether the HR
piControl goal of $-0.16$ is even the right target for this configuration should be revisited, given that
$\approx0.67$ of any coupled value is an atmosphere matching observations.

\subsection{Result II --- the boreal \emph{summer} cold bias is atmospheric}\label{sub:amipsummer}
Monthly 2\,m-T bias against CRUNCEP3 (1901--10), land only, cos-lat weighted:

\begin{center}\small
\begin{tabular}{l r r r}
\toprule
series & Boreal JJA/ann & Siberia JJA/ann & E.\ Siberia JJA/ann \\
\midrule
\textbf{AMIP (atmosphere only)} & $-1.64$ / $-1.05$ & $\mathbf{-2.16}$ / $-0.89$ & $\mathbf{-2.24}$ / $-0.49$ \\
080a (baseline, old sea ice)    & $-0.36$ / $-0.90$ & $-1.41$ / $-1.36$ & $-1.80$ / $-1.48$ \\
09A (baseline + newSeaIce)      & $-1.24$ / $-2.06$ & $-2.19$ / $-3.00$ & $-2.47$ / $-2.89$ \\
09B (06T + newSeaIce)           & $-1.56$ / $-2.33$ & $-2.50$ / $-3.11$ & $-2.72$ / $-3.04$ \\
09C (06V + newSeaIce)           & $-1.93$ / $-2.64$ & $-2.54$ / $-3.08$ & $-2.69$ / $-3.05$ \\
\bottomrule
\end{tabular}
\end{center}

The atmosphere alone reproduces essentially all of 09A's Siberian JJA bias ($-2.16$ vs $-2.19$) and most of
09B/09C's. \textbf{Removing the ocean and LPJ-GUESS does not remove the summer cold bias.} The cold-season
half behaves oppositely --- AMIP's Siberian \emph{annual} bias is only $-0.89$ against $-3.0$ coupled --- so
September--April remains the coupled/vegetation problem of Part~II, and JJA is now a separate, atmospheric one.

\subsubsection{Correction (2026--07--29): about half of that $-2.16$\,K is the reference period}\label{sub:perioffset}
The number above is scored against \textbf{ERA5 1990--2014}, while these runs use \textbf{1870s observed SST
with 1850 GHG}. Roughly $130$ years of greenhouse warming therefore sits inside a quantity the campaign has
been treating as model error --- and tuning against. Measured from the DKRZ ERA5 pool
(\texttt{scripts/analysis/era5\_period\_offset.sh}), Siberian JJA over the same box used for evaluation:

\begin{center}
\begin{tabular}{lr}
\toprule
ERA5 period & offset vs 1990--2014 \\
\midrule
1979--1989 & $\mathbf{-0.670}$\,K \\
1940--1969 & $\mathbf{-0.790}$\,K \\
\bottomrule
\end{tabular}
\end{center}

\noindent The interannual sd is $0.578$\,K, giving SE $=0.209$\,K on the 1979--89 vs 1990--2014 difference
($t\approx3.2$): the offset is real. ERA5 stops at 1940, so chaining through HadCRUT5's global series
(1870s $\to$ 1990--2014 $= +0.756$\,K) scaled by ERA5's own Siberian-JJA-to-global amplification over the
overlapping window ($\times1.48$) implies a Siberian JJA period offset of $\approx1.12$\,K, leaving a
\textbf{genuine model cold bias of $\approx-1.0$\,K rather than $-2.2$\,K}.

\textbf{This changes the target, not the diagnosis.} The attribution above is unaffected --- it is a
model-versus-model comparison (AMIP against coupled 09A) in which any reference offset cancels, so ``the JJA
bias is atmospheric'' stands. What changes is \emph{how much} of it we should try to remove: closing the full
$2.2$\,K would manufacture a $\approx1$\,K \emph{warm} bias in the coupled pre-industrial state.

\emph{Caveats.} The ERA5 box is unmasked whereas the model figures are land-masked; land warms faster than
ocean, so the true land-only offset is if anything larger and this correction is conservative. The $\times1.48$
amplification is extrapolated back from a single window, and HadCRUT5's 1870s global mean rests on sparse
coverage. \textbf{The same flaw applies to the CERES-based energy targets} (\S\ref{sub:ceres},
\S\ref{sub:sobands}): CERES EBAF here is the 07/2005--06/2015 climatology scored against an 1870s model, so
the $+8.08$\,W/m$^2$ SO SW CRE gap and the $+0.53$\,W/m$^2$ imbalance carry an analogous, \emph{as yet
unquantified}, period offset. The decisive fix for both is a single present-day AMIP leg (1989--2015,
transient GHG) scored directly against ERA5 1990--2014 and CERES 2005--2015; the AMIP SST forcing already
covers 187001--201512.

\textbf{Does this shift the residual onto LPJ-GUESS's CRUNCEP calibration?} Partly, but the corroboration
attempted on the \emph{temperature} axis fails: NCEP2 Siberian JJA is $-0.118$\,K \emph{colder} than ERA5 over
1990--2014 (paired by year, $t=-2.6$) --- real but negligible, and the wrong sign for a story in which
LPJ-GUESS grew its trees under an over-warm forcing. That is consistent with CRUNCEP3's construction:
temperature derives from CRU, bias-corrected against stations, while \emph{radiation} derives from NCEP, which
is precisely where the $+21$\,W/m$^2$ excess against CERES was measured (\S\ref{sub:ceres}). The calibration
mismatch sits on the radiation axis. The direct evidence remains the forcing-transfer test
(\S\ref{sub:xfer}), which does not depend on which variable carries it.

\begin{figure}[H]\centering
\includegraphics[width=0.99\textwidth]{plots/campaign_boreal_seasonal_bias.png}
\caption{Boreal 2\,m-T bias by month against CRUNCEP3, land-masked. The atmosphere-only AMIP curve (blue)
sits at or below the coupled runs through JJA (shaded) but far above them in winter: \textbf{the summer
bias is atmospheric, the cold-season bias is not}. \emph{Caveat:} the ``Boreal 55--70\textdegree N'' panel is
contaminated by the Greenland ice sheet (JJA snow depth $320$\,mm w.e.\ in that box vs $5.5$\,mm for
Siberia); the Siberia and E.\ Siberia panels are clean and carry the argument.}
\end{figure}

\textbf{Superseded numbers.} \dead{AMIP Siberia JJA $-1.23$, E.\ Siberia $-1.59$, i.e.\ indistinguishable from
the coupled baseline, so coupling adds nothing in summer.} Those came from the earlier
\texttt{amip\_pi\_forcing} run, built on the wrong OIFS branch and with \texttt{TMPRALBSEAD}$=0.075$. The
corrected build is $\approx0.9$\,K colder in Siberian JJA; the conclusion strengthens (the atmosphere is now
\emph{colder} in summer than the coupled baseline) but the numbers changed.

\section{Round 09 --- new sea-ice thermodynamics (2026--07--27/28)}\label{sec:round09}
FESOM 2.7 \texttt{awiesm3-implicit-ice-surftemp}\,@\,\texttt{db7dd9af} makes the ice
surface-temperature solve implicit (d$Q$/d$T$-linearised about the last transmitted \texttt{ist}) and hands
ownership of the ice surface energy balance to OpenIFS, which now runs a slab-conduction ice tile using
FESOM-supplied thickness; FESOM ice growth becomes $Q_\mathrm{atmice}=-a2ihf$. The pairing is verifiable at
runtime in the OASIS \texttt{namcouple}: the 09 runs carry a sixth o2a field,
\texttt{sit\_feom}$\to$\texttt{A\_Ice\_thickness}, absent from every 08 run.
Run identities and the 09B parameter block are catalogued in the 2026--07--28 update at the head of this
report; the quantitative comparison is below. \texttt{09C} was built and run on our own tree
(\texttt{Tuning\_test\_09C\_06V\_CRUNCEPinit\_newSeaIce}) to close the missing 06V corner.

\begin{center}\small
\begin{tabular}{l r r r r}
\toprule
run & net TOA [W/m$^2$] & global T2m [\textdegree C] & Siberia JJA [K] & Siberia ann [K] \\
\midrule
080a baseline, old sea ice & 1.442 & 13.606 & $-1.41$ & $-1.36$ \\
09A baseline + newSeaIce   & 1.286 & 13.148 & $-2.19$ & $-3.00$ \\
09B 06T + newSeaIce        & 1.171 & 12.794 & $-2.50$ & $-3.11$ \\
08B 06V, old sea ice       & 0.935 & 12.347 & $-3.11$ & $-4.09$ \\
\textbf{09C 06V + newSeaIce} & \textbf{0.919} & 12.492 & $\mathbf{-2.54}$ & $\mathbf{-3.08}$ \\
\bottomrule
\end{tabular}
\end{center}

On the clean old-vs-new 06V pair (08B\,$\to$\,09C) the radiation barely moves ($0.935\to0.919$) but the
boreal \textbf{warms by $+1.0$\,K in the Siberian annual mean}, $+0.57$\,K in JJA, $+0.15$\,K globally.
\textbf{09C dominates 08B on both objectives simultaneously} --- the first time in the campaign that has
happened, and the ``free warming'' sought in \S\ref{sec:rec}. But the effect is \emph{not} an offset: on the
baseline (080a\,$\to$\,09A) the same code change goes the other way, cooling $0.46$\,K globally and $1.6$\,K
in the Siberian annual mean. The sign depends on the ocean-mixing branch; no mechanism is established here,
and pinning it down needs the surface energy budget over ice.

\begin{figure}[H]\centering
\includegraphics[width=0.92\textwidth]{plots/campaign_net_toa_by_year.png}
\caption{Annual global-mean net TOA (thin) with 5-yr running mean (bold), common 1350--1379 window. The
dotted line is the \textbf{atmosphere-only floor} measured by AMIP ($+0.67$\,W/m$^2$); the shaded band is the
only part of the remaining gap that ocean-side tuning can still close. Everything above the dotted line is
atmospheric.}
\end{figure}

\emph{Bookkeeping caveat.} \texttt{sacct} shows a \texttt{Tuning\_test\_08A\_06T\_1hCPL\_MOSPP\_KPPLOW\_CRUNCEPinit}
completing all three chunks on 24--25 July, but no 08A directory exists on disk. 08A was the 06T-only
physical control --- Run~1 of the 2026-07-24 design note's own sequence --- so within round 08 every
surviving 06T-family entry (08C, 08E, 08G) carries an LPJ lever and the intended 06T-vs-06V comparison
cannot be made cleanly. Round 09 supersedes this, but conclusions drawn from round 08 alone inherit the hole.

\section{Why the summer bias exists --- a shortwave deficit (2026--07--28)}\label{sec:jja}
Siberia JJA surface energy budget, land-masked, cos-lat weighted (AMIP 1872--79 vs CRUNCEP3 1901--10, which
carries \texttt{rsns}/\texttt{rlns} and so permits a direct radiative comparison):

\begin{center}\small
\begin{tabular}{l r r r}
\toprule
 & AMIP & CRUNCEP3 & $\Delta$ \\
\midrule
2\,m temperature      & $9.74$\,\textdegree C & $11.92$\,\textdegree C & $\mathbf{-2.18}$\,K \\
SW net at surface     & $155.1$ & $188.3$ & $\mathbf{-33.2}$\,W/m$^2$ \\
LW net at surface     & $-49.2$ & $-64.5$ & $+15.2$\,W/m$^2$ \\
\midrule
net radiation         &         &         & $\mathbf{-18}$\,W/m$^2$ \\
\bottomrule
\end{tabular}
\end{center}

Supporting AMIP diagnostics: SW$\downarrow$ $188.6$, surface albedo $0.178$ (\texttt{fal} $0.172$),
LW$\downarrow$ $315.6$, SH/LH $-24.6/-56.8$ (\textbf{Bowen $0.43$}), cloud \texttt{tcc} $0.778$ /
\texttt{lcc} $0.505$ / \texttt{mcc} $0.386$, \texttt{skt}$-$\texttt{2t} $+0.43$\,K, \texttt{swvl1} $0.343$,
snow depth $5.5$\,mm w.e.

The LW term is \emph{compensation} (a colder, cloudier surface loses less longwave), not cause. It is
\emph{not} albedo: a realistic boreal-summer $0.13$--$0.15$ instead of the model's $0.178$ buys only
$5$--$9$\,W/m$^2$. So SW$\downarrow$ itself is too low. This is a different mechanism from the winter
stable-BL inversion of \S\ref{sub:residual}, and it is not addressed by any lever proposed so far.

\subsection{Re-measured against CERES: cloud, not clear-sky --- and three times smaller}\label{sub:ceres}
CRUNCEP3 is a reanalysis-derived forcing product, so the deficit was re-measured against
\textbf{CERES EBAF Ed4.1} (satellite), which supplies all-sky \emph{and} clear-sky net SW at the surface and
so permits the decomposition
$\Delta(\mathrm{SW}_\mathrm{net})=\Delta(\mathrm{SW}_\mathrm{net,clear})+\Delta(\mathrm{SW\ CRE}_\mathrm{sfc})$
with no assumption about surface albedo or the direct/diffuse split.

\begin{center}\small
\begin{tabular}{l r r r r r r}
\toprule
JJA & \multicolumn{3}{c}{Siberia land} & \multicolumn{3}{c}{NH 55--75\textdegree N ocean} \\
\cmidrule(lr){2-4}\cmidrule(lr){5-7}
 & model & CERES & diff & model & CERES & diff \\
\midrule
SW net at surface, all-sky   & 155.1 & 166.3 & $\mathbf{-11.2}$ & 142.4 & 155.4 & $-13.0$ \\
SW net at surface, clear-sky & 246.6 & 241.4 & $+5.3$           & 241.9 & 252.6 & $-10.7$ \\
SW CRE at surface            & $-91.6$ & $-75.1$ & $\mathbf{-16.5}$ & $-99.5$ & $-97.3$ & $\mathbf{-2.2}$ \\
cloud area [\%]              & 77.8 & 69.6 & $\mathbf{+8.2}$   & 84.6 & 83.0 & $\mathbf{+1.6}$ \\
\bottomrule
\end{tabular}
\end{center}

Three results, all of which revise \S\ref{sec:jja}:

\textbf{(i) It is unambiguously a cloud error.} The decomposition closes exactly
($-11.2=+5.3-16.5$), and clear-sky surface SW is if anything slightly \emph{too bright}. Aerosol, water
vapour and surface albedo are therefore all exonerated --- and since PI aerosol is lower than present-day, the
PI-vs-CERES period mismatch works \emph{against} this finding, making it conservative.

\textbf{(ii) The deficit is $11$\,W/m$^2$, not $33$.}
\dead{Surface net SW over Siberian land in JJA is 33\,W/m$^2$ too low.} That figure was measured against
CRUNCEP3, which is itself $\approx21$\,W/m$^2$ brighter than satellite --- \textbf{CRUNCEP3 was the outlier,
not the model}. An $11$\,W/m$^2$ summer deficit remains the right order to explain a $\approx2$\,K land cold
bias, so the mechanism survives; only its magnitude changes.

\emph{Verified independently}, because the original comparison had assembled the two sides on different grids
with different land masks --- a way to manufacture exactly this kind of difference. Redone with all three
sources on the CERES $1^\circ$ grid under a single land mask (\texttt{verify\_ceres\_cruncep.py}): CERES
$166.7$, AMIP $155.8$ ($-10.8$), CRUNCEP3 $187.6$ ($+21.0$). Raising the land-mask threshold $0.5\to0.8$
changes these by $<1$\,W/m$^2$, and CERES's own JJA-mean file gives $168.4$ against $166.7$ from the monthly
climatology. The grid inconsistency was therefore not the cause.

\emph{Why CRUNCEP3 is bright} splits cleanly by season: in JJA it is the downward flux
($234$\,W/m$^2$ in July against CERES's $209$, the NCEP high-latitude brightness bias), whereas in spring it
is the forcing file's fixed albedo (\S\ref{sub:ruledout}). CRUNCEP3 blends CRU \emph{station} temperature with
NCEP \emph{reanalysis} radiation, which is the natural reading of why its T2m is usable while its SW is not.

\begin{figure}[H]\centering
\includegraphics[width=0.99\textwidth]{plots/verify_ceres_amip_cruncep.png}
\caption{Surface net shortwave against CERES EBAF, all sources on a common $1^\circ$ grid with one land mask.
AMIP (orange) is $\approx11$\,W/m$^2$ \emph{dark} in JJA; the CRUNCEP forcing (green) is $\approx21$\,W/m$^2$
\emph{bright} in JJA and far more so in spring, where its fixed albedo does not mask snow. The two forcings
differ from each other by $\approx32$\,W/m$^2$ in the growing season --- the transfer step of
\S\ref{sec:calibmismatch}.}
\end{figure}

\textbf{(iii) It is specific to boreal land.} Over ocean in the same latitude band the cloud error nearly
vanishes ($-2.2$\,W/m$^2$ CRE, $+1.6$\,pp area) while over Siberian land it is $-16.5$ and $+8.2$. A
\emph{global} cloud knob would therefore move the ocean, where nothing is wrong, and spend TOA budget for no
benefit. The error is over land, in summer --- pointing at surface-driven convection and boundary-layer cloud
rather than the large-scale cloud scheme.

\textbf{Amount or opacity?} Both, roughly half each. At the observed CRE per unit cloud area
($-75.1/0.696\approx-108$\,W/m$^2$), the $+8.2$\,pp excess accounts for $\approx-8.8$ of the $-16.5$\,W/m$^2$
CRE error; the residual $\approx-7.7$ is opacity or vertical placement. A cloud-\emph{amount} lever alone
therefore reaches at most half of it.

\subsection{Does the cloud fix conflict with the energy target? Less than first thought}\label{sub:jjaknob}
\dead{Burning off that cloud stops reflecting it at TOA, costing $\approx+0.18$\,W/m$^2$ on the annual global
TOA --- a $27\,\%$ worsening --- so a pure cloud fix cannot meet a fixed-TOA constraint, and only a
radiatively neutral surface-flux lever can.} Two things falsify that framing:

\begin{itemize}\itemsep1pt
\item \textbf{The target is the \emph{surface}, not TOA} (\S\ref{sub:snowenth}), and the drift it controls is
driven by the flux into the \emph{ocean}. The cloud error is over \emph{land}
(\S\ref{sub:ceres}), where annual-mean energy storage is $\approx0$ --- land returns what it receives. A
boreal-land cloud correction therefore barely touches the quantity that governs spin-up drift.
\item \textbf{The error is land-specific}, so the correction need not be a global knob at all
(\S\ref{sub:ceres}) --- which was the whole reason a TOA cost was incurred in the first estimate.
\end{itemize}
The two objectives are consequently far more separable than the original framing implied, and the cloud
error can be attacked at its diagnosed cause rather than through a proxy.

\subsection{The stomatal-resistance sensitivity test (experiment A)}\label{sub:expA}
A second, radiatively neutral lever was nonetheless identified and is being tested, because the measured
Bowen ratio over Siberian JJA land is $\mathrm{SH/LH}=-24.6/-56.8=\mathbf{0.43}$, low against the general
observed picture for boreal conifer forest. Minimum canopy resistance for boreal needleleaf
(\texttt{susveg\_mod.F90}:266--267) is $250$\,s/m, with $500$ retained just above as a comment.

\textbf{Provenance --- and a correction.} \dead{$500$ is the IFS default and has been halved in this build.}
Git refutes this: both the active $250$ and the commented $500$ are present in
commit \texttt{cf82033}, ``initial commit of 48r1 released version'', and no AWI commit has ever touched
\texttt{RVRSMIN}. \textbf{$250$ is what ECMWF ships in 48r1}; $500$ is the older HTESSEL-era value they
themselves revised. Setting $500$ moves \emph{away} from the released configuration, so this is a
\textbf{sensitivity test, not a default restoration}.

\emph{Status:} \texttt{amip\_expA\_rvrsmin500}, 6\,yr, paired against \texttt{amip\_pi\_base} on 1872--1875.
Verified compiled in by inspecting the rebuilt object file (the value appears twice for types 3--4; $250$
survives only for Desert and Mixed Forest). Note this tests \emph{partitioning}, not the diagnosed cause: if
the cloud error is real then a resistance change compensates for it, and with $11$\,W/m$^2$ less available
energy there is also less to repartition. \emph{Caveat:} ``Bowen $0.43$ is too low'' is inferred from the
literature's general picture, not measured --- neither CRUNCEP3 nor CERES supplies SH/LH here, and FLUXNET
boreal sites would settle it.

\textbf{A connection worth testing in the same run.} Excess summer evapotranspiration over boreal land would
\emph{both} cool the surface directly \emph{and} moisten the boundary layer, feeding the excess low cloud
(model \texttt{lcc}$=50.5\,\%$ over Siberian land in JJA). If cutting transpiration also reduces low cloud,
the partitioning error and the cloud error are the same problem, and experiment A tests both at once.

\section{Forcing transfer: the spin-up was grown under brighter, warmer conditions than the model delivers}\label{sec:calibmismatch}
The vegetation state every coupled run in rounds 08 and 09 restarts from is the standalone offline spin-up
\texttt{LPJG-SPINUP\_2000Y\_TCO95\_CORE3\_CRUNCEPcalibrated\_v3}
(\texttt{/work/bb1469/a270270/runtime/lpjg-spinup/}), driven by
\texttt{CRUNCEP\_noLPJG\_1d\_1901-1910\_TCO95\_calibrated\_v3.nc} (verified at
\texttt{lpj\_guess\_levante\_spinup\_...\_2000y.yaml}:38). \textbf{That forcing and the atmosphere the coupled
model provides differ substantially, and in the same direction for both light and temperature.}

\begin{center}\small
\begin{tabular}{l r r}
\toprule
Siberia land, JJA & surface net SW [W/m$^2$] & vs CERES \\
\midrule
CERES EBAF (observed)                       & 166.7 & --- \\
\textbf{CRUNCEP forcing} (what LPJG grew under) & \textbf{187.6} & $\mathbf{+21}$ ($+13\,\%$) \\
\textbf{AMIP} (what our atmosphere delivers)    & \textbf{155.8} & $\mathbf{-11}$ ($-7\,\%$) \\
\midrule
\multicolumn{2}{l}{\textbf{gap between the two forcings}} & $\mathbf{32}$\,W/m$^2$ ($\approx20\,\%$) \\
\multicolumn{2}{l}{\textbf{plus 2\,m temperature}} & $\mathbf{2.2}$\,K (11.9 vs 9.7\,\textdegree C) \\
\bottomrule
\end{tabular}
\end{center}

The errors point \emph{opposite} ways and therefore compound: CRUNCEP is too bright (the NCEP high-latitude
brightness bias --- its downward SW is $234$\,W/m$^2$ in July against CERES's $209$, $\approx12\,\%$ high),
while our atmosphere is too dark by the cloud error of \S\ref{sub:ceres}. Vegetation equilibrated to the
first and then handed the second experiences a $\approx20\,\%$ drop in growing-season radiation and a
$2.2$\,K cooling --- a step change in exactly the direction that pushes a boreal system from trees toward
grasses. All numbers here are computed on a common CERES $1^\circ$ grid with a single land mask
(\texttt{verify\_ceres\_cruncep.py}); mask threshold $0.5\to0.8$ moves them by $<1$\,W/m$^2$.

\subsection{Ruled out by measurement}\label{sub:ruledout}
Several components were suspected along the way and cleared. Recording them so the same ground is not
re-covered:

\begin{center}\small
\begin{tabular}{l l}
\toprule
suspect & the observation that cleared it \\
\midrule
LPJ-GUESS core radiation handling &
\dead{PAR is inflated because no albedo is applied.} Reading \texttt{rsns} as
\texttt{NETSWRAD\_TS} and applying \emph{no} further albedo is \textbf{correct} for a net-flux \\
& input (\texttt{driver.cpp}:1435--1449); a second correction would double-count. Model behaves as designed. \\[2pt]
EC-Earth spin-up interface &
Reads the right variable and sets the right type (\texttt{framework.cpp}:2501--2507). No fault found. \\[2pt]
The forcing script's fixed albedo &
\dead{\texttt{rsns}$=$\texttt{FSDS}$\times0.861$ is a fixed 0.139 albedo, so spring PAR is inflated $\times1.66$
and the spin-up is undermined.} \\
& True that the albedo is fixed (ratio $0.861$ in every cell, April and July alike). But in JJA that gives \\
& $0.139$ against CERES's $\mathbf{0.142}$ --- right to $0.003$. The large spring error
($0.139$ vs $0.535$ in April) \\
& is \textbf{biologically inert}: measured over 687 Siberian cells at equilibrium, April carries
$\mathbf{1.1\,\%}$ of annual \\
& GPP and April LAI is $0.02$ (no canopy), while \textbf{JJA carries $\mathbf{80\,\%}$}. Weighted by when carbon \\
& is actually assimilated, the PAR excess is $\approx13\,\%$, not $66\,\%$. \\[2pt]
\dead{...and it is an unusual, campaign-specific defect} &
Standalone LPJ-GUESS assumes a constant albedo natively: \texttt{BETA}$=0.17$
(\texttt{driver.cpp}:1320), applied \\
& in both the sunshine and \texttt{SWRAD} paths. The script's $0.139$ is a slightly darker version of \\
& standard practice, and switching to its \texttt{bulk\_albedo} option would give $0.195$ in JJA --- \\
& \emph{worse} than the fixed value. \textbf{Do not ``fix'' the script; it is not the problem.} \\
\bottomrule
\end{tabular}
\end{center}

\subsection{What survives, and how to settle it}\label{sub:forcingtest}
What remains is not a defect in any component but a \textbf{transfer problem}: a $\approx20\,\%$
growing-season radiation step and a $2.2$\,K temperature step between the forcing the vegetation
equilibrated to and the forcing it is then asked to live in. \textbf{Whether LPJG's \emph{parameters} were
tuned against CRUNCEP --- the ``calibrated'' in the directory name --- is not established}; the 07-series did
change \texttt{greff\_min} and \texttt{pstemp\_low}, but no evidence was found of what was calibrated against
what, and the folder name is not proof.

\textbf{It still weakens the ``competition, not climate'' verdict of \S\ref{sub:decisive}}, though less than
first argued. That test moved a standalone re-spin from the CRUNCEP forcing to the OIFS--AMIP forcing and
read the resulting under-vegetation as an LPJG-internal competition failure. But that move alone removes
$\approx20\,\%$ of the growing-season light and $2.2$\,K, so under-vegetation is partly what the forcing
change predicts on its own. The verdict should be read as provisional. (The forcing used there was also the
\emph{superseded} AMIP build, \S\ref{sub:amipsummer}.)

\subsection{The transfer test, measured: forcing alone halves boreal tree cover}\label{sub:xfer}
The test needs no knowledge of what was tuned --- take the \emph{same} LPJG parameters and spin to
equilibrium under each forcing; the difference \emph{is} the transfer penalty. \textbf{Both arms already
existed} and required no new computation:
\begin{center}\footnotesize
\begin{tabular}{l l}
\toprule
CRUNCEP arm & \texttt{lpjg-spinup/LPJG-SPINUP\_2000Y\_TCO95\_CORE3\_CRUNCEPcalibrated\_v3} \\
AMIP arm    & \texttt{lpjg-spinup-develop/LR\_2000y\_PIforcing\_v3} \\
\bottomrule
\end{tabular}
\end{center}
Controlled: \texttt{global.ins} \textbf{byte-identical} between the two trees, \texttt{firemodel "GLOBFIRM"}
in both \texttt{.ins} files, \texttt{nyear\_spinup}$=2000$ / \texttt{freenyears}$=100$, fixed 1850
CO$_2$/Ndep/LU, both reaching year 3900, and \textbf{10\,043 gridcells in common} (99.8\,\% of the AMIP
gridlist). The material difference is the forcing file.

\begin{center}\small
\begin{tabular}{l l r r r r}
\toprule
region & & CRUNCEP-forced & AMIP-forced & diff & \% \\
\midrule
\textbf{Siberia} & \texttt{TREEFPC}  & 0.357 & \textbf{0.163} & $-0.194$ & $\mathbf{-54}$ \\
                 & \texttt{GRASSFPC} & 0.288 & 0.330 & $+0.042$ & $+15$ \\
                 & \texttt{AGDD5}    & 776 & 645 & $-131$ & $-17$ \\
\midrule
\textbf{E.\ Siberia} & \texttt{TREEFPC}  & 0.359 & \textbf{0.122} & $-0.236$ & $\mathbf{-66}$ \\
                     & \texttt{GRASSFPC} & 0.316 & 0.398 & $+0.082$ & $+26$ \\
                     & \texttt{AGDD5}    & 748 & 568 & $-180$ & $-24$ \\
\midrule
NH 45\textdegree N$+$ & \texttt{TREEFPC} & 0.383 & 0.287 & $-0.096$ & $-25$ \\
\bottomrule
\end{tabular}
\end{center}

\textbf{Changing nothing but the forcing halves the Siberian tree fraction and cuts East Siberia's by
two-thirds, with grass taking up the space.} For scale, the coupled 080a run lost \texttt{TREEFPC}
$0.310\to0.121$ over Siberia in 50 years (\S\ref{sub:cruncep3_init_followup}); the forcing swap alone
reproduces essentially that entire loss at equilibrium, with no parameter changed. \texttt{AGDD5} falls
$17$--$24\,\%$ in the same direction, so this is not a light effect alone --- the $2.2$\,K also does real work
through the growing-degree-day gates.

\begin{figure}[H]\centering
\includegraphics[width=0.99\textwidth]{plots/forcing_transfer_test.png}
\caption{Zonal means at equilibrium (yr 3900), 10\,043 common gridcells, byte-identical \texttt{global.ins}.
South of $\approx52$\textdegree N the two arms are indistinguishable; northward they diverge sharply and the
AMIP-forced treeline retreats several degrees south. The deficit the campaign has been attributing to
competition is reproduced by the forcing change alone.}
\end{figure}

\textbf{Three caveats, all material.}
\begin{enumerate}\itemsep1pt
\item The arms were built from \emph{different binaries} (a270092 vs a270270) of the same LPJG version with
identical \texttt{global.ins}. A compiled difference cannot be excluded; re-running both arms from one build
($\approx3$\,h) would close it, and should be done before this is published.
\item The AMIP forcing is the \emph{superseded} build (\S\ref{sub:amipsummer}), $\approx0.9$\,K warmer over
Siberian JJA land than the corrected one --- so this is a \textbf{lower bound} on the penalty against today's
atmosphere.
\item \dead{\texttt{FORESTFPC}$=0$ in both arms, so the absence of closed canopy survives as a separate
unexplained problem.} \texttt{FORESTFPC} is a \emph{void diagnostic} in this configuration --- it counts only
the managed-forest land-use class, which is empty here, so it is identically zero regardless of what the
vegetation does. There is no separate closed-canopy failure to explain.
\end{enumerate}

\section{Revised plan: tune the atmosphere on AMIP}\label{sec:amipplan}
Both errors that remain --- $\approx0.83$\,W/m$^2$ of TOA imbalance (\S\ref{sub:amiptoa}) and
$\approx2.2$\,K of Siberian summer cold (\S\ref{sub:amipsummer}) --- are atmospheric, and both are now
\emph{measured} in a configuration where neither the ocean nor LPJ-GUESS can confound the answer. They should
therefore be tuned on AMIP, not in coupled mode. This does not retire Phases~1A/1B of \S\ref{sec:workplan}
--- the LPJG competition fix and the stable-BL mixing fix still target the \emph{cold-season} bias --- but it
inserts a cheaper, better-posed stage ahead of Phase~2.

\textbf{Two objectives, now known to be largely separable} (\S\ref{sub:jjaknob}):
\begin{enumerate}\itemsep1pt
\item \textbf{Surface energy balance $\to0$} to cut spin-up drift. This is an \emph{ocean} surface-flux
problem; the relevant number is $F_\mathrm{sfc}$ over ocean, which has \emph{not} yet been computed --- all
figures quoted so far ($+0.52$ AMIP, $+0.95$ 09C, $+1.77$ 080a) are global. \textbf{Open item.}
\item \textbf{Boreal summer warmth} to restore GDD5 and let trees compete. This is a \emph{land} problem,
caused by an $11$\,W/m$^2$ cloud-driven SW deficit that is absent over ocean (\S\ref{sub:ceres}).
\end{enumerate}

\begin{center}\small
\begin{tabular}{l l l l}
\toprule
 & change & where & status / expectation \\
\midrule
\textbf{A} & \texttt{RVRSMIN(3,4)}: $250\to500$ & \texttt{susveg\_mod.F90}:266--267 & \emph{running}; radiatively
neutral; tests partitioning, not the diagnosed cause \\
B & \texttt{RVLAI(3,4)}: $5.0\to4.0$ & \texttt{susveg\_mod.F90}:133--134 & same sign, smaller \\
C & \dead{\texttt{RCLDIFF} $\uparrow$, global} & \dead{\texttt{NAMCLDP}} & \dead{$+$surface SW, but $+$TOA} \\
C$'$ & a \emph{land}-acting cloud lever & convection / BL cloud & the error is land-only, so a global knob
would move the ocean where nothing is wrong \\
D & A\,$+$\,C$'$ & & joint target \\
\bottomrule
\end{tabular}
\end{center}

\textbf{Gates:} Siberian JJA bias toward $0$ \emph{and} SW CRE toward CERES; ocean $F_\mathrm{sfc}$ not
degraded; no new error outside the boreal (judge on the global \texttt{part8}/CMPI pattern, \S\ref{sec:cmpi}).
Since the cloud error splits roughly half amount, half opacity (\S\ref{sub:ceres}), a lever acting only on
cloud amount can close at most half of it, and the diagnostic to watch is SW CRE against CERES rather than
temperature alone.

\textbf{Why AMIP is the right testbed.} $2.3$ node-hours per model year against $15.8$ coupled --- a
$\approx7\times$ saving --- and 10\,yr completes in one job. Measured interannual scatter of the boreal JJA
mean is $\sigma=0.25$\,K (Boreal), $0.59$\,K (Siberia), $0.72$\,K (E.\ Siberia), so with 8 usable years the
standard error is $0.09$--$0.26$\,K against biases of $1$--$2$\,K: sampling is not the binding constraint.
\emph{Land spin-up is} --- with \texttt{lresume:\ false} the deep soil (\texttt{stl4}, \texttt{swvl4}) carries
initialisation shock, so discard the first 1--2 years. For paired tuning experiments 5\,yr (1 discarded
$+$ 4 used) suffices, since identical prescribed SST makes most interannual variance cancel in the difference.

\clearpage
\section{Round 10 design: a two-pronged, targeted attack (2026--07--28)}\label{sec:round10}
\emph{Original design as proposed. Operator notes to be appended in \S\ref{sub:r10notes}.}

\textbf{The structural fact the design rests on.} The two largest regional radiation errors have
\emph{opposite} sign at TOA (\S\ref{sub:sobands}, \S\ref{sub:ceres}), so fixing them together is
self-balancing, whereas fixing either alone is not:

\begin{center}\small
\begin{tabular}{l l l r}
\toprule
 & error & fixing it & $\Delta$TOA \\
\midrule
\textbf{A1 Southern Ocean} & too \emph{little} cloud ($-6.6$\,pp, SW CRE $+7.8$) & adds cloud, reflects more &
$\mathbf{-0.5}$ to $-0.6$ \\
\textbf{A2 boreal land} & too \emph{much} cloud ($+8.2$\,pp, surface SW $-11$ in JJA) & removes cloud, absorbs more &
$\mathbf{+0.1}$ \\
\bottomrule
\end{tabular}
\end{center}

\subsection{Prong A --- AMIP (atmosphere), A1 before A2}
\textbf{A1 first.} It moves the global target in the right direction ($+0.67\to\approx+0.1$) \emph{and}
attacks the error the campaign has suspected since \S\ref{sec:so}. A2 then costs back only $\approx0.1$ of the
$\approx0.55$ that A1 buys, so the boreal fix becomes affordable --- the ``warm floor'' logic of Part~III,
relocated from the vegetation to the atmosphere.

\textbf{Critical: in AMIP the Southern-Ocean target is \emph{radiation only}, never temperature.}
SST and sea ice are prescribed, so the SO surface temperature cannot respond to a radiation error --- the
excess absorbed shortwave is simply exported through the fixed boundary. \textbf{A1 must therefore be judged
on SW CRE and cloud area against CERES, and it is meaningless to look for an SO temperature response in
AMIP.} This is also why the SO \emph{warm} bias appears only in AWI-CM3/AWI-ESM3: it is a \emph{coupled}
phenomenon, in which the same $+5.8$\,W/m$^2$ of excess absorbed energy has somewhere to go and warms the
ocean. The causal chain is then explicit and closes the compensating-error argument of \S\ref{sec:so}:
\begin{center}\small
AMIP SO cloud deficit ($-6.6$\,pp) $\to$ $+5.8$\,W/m$^2$ excess absorbed $\to$
\emph{(coupled only)} SO warm SST bias $\to$ suppressed by the 06T/06V ocean-mixing levers
\end{center}
i.e.\ Part~I was damping Weddell/Ross convection to compensate an \emph{atmospheric} cloud error. It follows
that the only test of whether A1 fixed the real problem is a \emph{coupled} run --- AMIP can confirm the
radiation is right, but not that the SO warm bias is gone.

\textbf{On levers.} OpenIFS has no land-only cloud knob, and that is informative rather than obstructive: the
land-specificity of A2 ($-16.5$\,W/m$^2$ CRE error over Siberian land against $-2.2$ over ocean at the same
latitude) is itself evidence that the cause is \textbf{land-surface moisture supply}, not the cloud scheme ---
which makes the \texttt{RVRSMIN} sensitivity test (\S\ref{sub:expA}) potentially the actual fix rather than a
proxy for one. A1, being over ocean, is a genuine cloud-scheme / microphysics target.

\textbf{Guardrails on every A-run} --- all four, every time:
\begin{enumerate}\itemsep1pt
\item global net TOA \emph{and} surface flux (state which convention, \S\ref{sub:snowenth});
\item SO 45--65\textdegree S SW CRE and cloud area against CERES --- \emph{radiation only}, since prescribed
SST makes an SO temperature response impossible in AMIP by construction (above);
\item \textbf{the tropics} --- $50\,\%$ of the globe and the largest single contribution ($-1.26$), and the
term the 06O/06T/06V levers degraded (\S\ref{sec:cmpi});
\item CMPI by \emph{pattern}, not the global number.
\end{enumerate}

\subsection{Prong B --- LPJ-GUESS, and it must go second}
\textbf{Do not tune competition parameters yet.} The tree deficit is largely forcing-driven ($-54\,\%$
Siberian \texttt{TREEFPC} from the forcing swap alone, \S\ref{sub:xfer}), so tuning LPJG while its forcing is
$11$\,W/m$^2$ dim and $2.2$\,K cold would bake the atmospheric error into vegetation parameters. That is the
compensating-error trap, and it is what the 07-series was already doing --- which is the natural reading of
why it achieved nothing (\S\ref{sec:round09}).

\begin{enumerate}\itemsep1pt
\item Complete Prong A.
\item \textbf{Regenerate the LPJG forcing} from the corrected atmosphere. Enabling work required:
\texttt{amip\_pi\_base} uses the default \texttt{file\_def} and emits the evaluation fields but \emph{not} the
daily LPJG forcing set, so a merged \texttt{file\_def} (or a second short run) is needed.
\item \textbf{Re-spin and re-derive the residual.} Whatever tree deficit survives against a corrected forcing
is the real LPJG target. Only then are \texttt{greff\_min} / \texttt{pstemp\_low} legitimate levers.
\end{enumerate}

\textbf{Gates, replacing the void one.} \dead{FORESTFPC$>0$} is unachievable by construction
(\S\ref{sub:xfer}) and must not be used. Use instead: \texttt{TREEFPC} against the CRUNCEP-arm equilibrium
($0.357$ Siberia, $0.359$ E.\ Siberia); \texttt{AGDD5} recovery toward $776$/$748$; and \texttt{cvh} toward
the MODIS $0.4$--$0.7$ closed-taiga range.

\subsection{Ordering}
\textbf{A1 $\to$ A2 $\to$ regenerate forcing $\to$ re-spin LPJG $\to$ B.} Prong B cannot be evaluated honestly
until Prong A has fixed what feeds it.

\subsection{Operator notes}\label{sub:r10notes}
\emph{(reserved --- to be filled in)}

\appendix
\section{Parameter catalogue (distinguishing change per run)}\label{app:params}
\begin{center}\footnotesize
\begin{tabular}{l l l l}
\toprule
Run & Change (file/namelist) & Baseline\,$\to$\,New & Process \\
\midrule
06A & \texttt{albpnd} (\texttt{namelist.ice}) & $0.2\to0.28$ & melt-pond albedo \\
06B & +\,\texttt{rfracmax} & $1.0\to0.75$ & pond water retention \\
06C & +\,\texttt{pndaspect} & $0.8\to1.3$ & pond depth:area (=common pond) \\
06D \emph{HRlike} & (only common pond shift) & -- & \textcolor{red}{no other diff found} \\
06G & \texttt{albpnd/rfracmax/pndaspect} & $0.35/0.6/1.6$ & stronger meltpond \\
06L & \texttt{albpnd/rfracmax/pndaspect} & $0.4/0.55/1.8$ & extra-strong meltpond \\
06H & strong pond + \texttt{RVICE} (\texttt{\&NAMCLDP}) & $0.16\to0.18$ & meltpond + cloud-ice fall speed \\
06Q & \texttt{albw/albocn} & $0.1\to0.045$ & open-water/ocean albedo \\
06R & \texttt{albw/albocn} & $0.1\to0.075$ & open-water/ocean albedo \\
06S & \texttt{whichevp} (\texttt{\&ice\_dyn}) & $1\to0$ & sea-ice rheology mEVP$\to$EVP \\
06E & \texttt{RVICE}+\texttt{GGAUSSB} (\texttt{\&NAMGWWMS}) & $0.18$; $-0.5\to-0.6$ & cloud-ice + gravity-wave \\
06F & \texttt{LRDALB} (\texttt{NALBEDOSCHEME=3}) & F\,$\to$\,T & diagnosed surface albedo \\
06P & \texttt{ENTSTPC3} (\texttt{\&NAMCUMF}) & absent\,$\to$\,1 & deep-convection entrainment \\
06N & \texttt{use\_momix} (\texttt{namelist.tra}) & F\,$\to$\,T & near-surface ocean mixing (mospp) \\
06M & OASIS coupling period + LAG & $7200\to3600$\,s & atmos$\leftrightarrow$ocean 2h$\to$1h (1hcpl) \\
06O & 1hcpl + mospp & -- & coupling + mixing \\
06O4 & O + \texttt{momix\_kv} & $0.01\to0.012$ & + larger M-O mixing \\
06O5 & O + \texttt{momix\_kv} & $0.01\to0.02$ & + even larger mixing \\
06T & O + KPP-low (\texttt{kpp\_av0/kv0} $0.005\to0.003$, bckg $\downarrow10\times$) & -- & + weaker KPP mixing \\
06U & T + open-water albedo $0.075$ & -- & combo \\
06V & T + \texttt{ENTSTPC3=1} & -- & combo (most aggressive) \\
\bottomrule
\end{tabular}
\end{center}

\section{Implementation \& lab-notebook detail}\label{app:impl}
Operational detail moved here to keep the work-plan science legible. Retained for reproducibility.

\textbf{A.1 esm-tools fixes for the AMIP--noLPJG generator (\S\ref{sub:amipgen}).} Model tree \texttt{oifsamip-cy48}
git-aligned to \texttt{awiesm3-develop} (oifs \texttt{local\_combined\_fixes}\,@\,1303a90, oasis \texttt{afac6a3},
xios\,\texttt{2.5.3}) + bundle/field patches. Three bugs fixed, committed \& pushed (commit \texttt{9921d95b}, branch
\texttt{development\_interactiv\_mesh\_awiesm3}, and a follow-up for the WAM recipe):
\begin{itemize}\itemsep0pt
\item \texttt{xios.yaml} dereferenced the \texttt{fesom} namespace via \texttt{choose\_fesom.*} and crashed \emph{any}
FESOM-less setup --- now gated on \texttt{"\$((\;'fesom' in config\;))"}.
\item \texttt{oifsamip.yaml} never overrode the OIFS \texttt{comp\_command}, so the binary compiled \emph{without} XIOS
or coupling despite \texttt{LXIOS:TRUE} --- now adds \texttt{--with-cplng2 --with-oifs-xios --build-type=release}
(like awicm3/awiesm3).
\item WAM aborted (\texttt{CDATEA}$\neq$\texttt{CDATEF}) at any start year $\neq$1990 because the OIFS
\texttt{preprocess} step (\texttt{change\_icm\_date.sh}, which re-dates ICMGG/ICMSH/wave init files) never ran:
\texttt{oifsamip.yaml} lacked \texttt{prepcompute\_recipe:\;\$\{oifs.prepcompute\_recipe\}} and fell back to the
default recipe with no preprocess step. Added (matches \texttt{awicm3.yaml}).
\end{itemize}
Runscript \texttt{oifsamip-cy48-levante-TCO95L91\_lpjgforcing.yaml} overrides XIOS output to
\texttt{file\_def\_lpjg\_spinup} + \texttt{field\_def\_lpjg\_safe}, emitting daily
\texttt{tas,tasmin,tasmax,pr,rsns,rlns,sfcWind} + \texttt{tdps,ps,psl,clt,prsn,evspsbl,prw,sd,rsn}.
\emph{esm\_master gotcha:} the single-component target is \texttt{recomp-<setup>/oifs-48r1} (the \emph{component}
name); the alias \texttt{/oifs} silently falls back to the setup and doubles the path
(\texttt{oifsamip-cy48/oifsamip-cy48}).

\textbf{A.2 Flux-accumulation fix (\S\ref{sub:amipgen}), at root.} The de-accumulation divisor in
\texttt{field\_def\_lpjg\_safe.xml} was the hard-coded \texttt{/21600} (6\,h). It is now templated:
\texttt{field\_def\_lpjg\_safe.xml.j2} uses \texttt{/\{\{oifs.flux\_accumulation\_period\}\}}, with
\texttt{oifs.flux\_accumulation\_period} defaulting to \texttt{\$\{oifs.time\_step\}} ($=$\,NFRHIS$\times$TSTEP; with
NFRHIS$=$1 this is the timestep). A validated $\times6$ scalar patch
(\texttt{AMIP\_noLPJG\_1d\_1870-1879\_TCO95\_PI\_fluxfix.nc}) corrected the existing forcing without a rerun.

\textbf{A.3 Phase-1B OpenIFS source edits (\S\ref{sub:phase1b}).} Paths under
\texttt{oifs-48r1/ifs-source/arpifs/phys\_ec/}:
\begin{center}\small
\begin{tabular}{l l l l l}
\toprule
knob & file:line & current & proposed & effect \\
\midrule
\texttt{ZLMIN} & \texttt{vdfexcu.F90}:201 & $30$\,m & $40$--$50$\,m & min stable mixing length \\
\texttt{ZCB}   & \texttt{vdfexcu.F90}:190 & $5.0$ & $3.0$--$4.0$ & LTG $b$: lower $\Rightarrow$ longer tail $\Rightarrow$ more $K$ \\
\bottomrule
\end{tabular}
\end{center}
Secondary companions if needed: the $0.1$ / $300$\,m stable mixing-length limits (\texttt{vdfexcu.F90}:374,376) and the
surface-layer constants \texttt{RCHETA}/\texttt{RCHETB}/\texttt{RCHBHDL} (\texttt{surf/module/suvexcs\_mod.F90}, mirror
in \texttt{suvdfs.F90}).

\section*{References}
{\footnotesize\setlength{\parindent}{0pt}\setlength{\parskip}{2.5pt}
Alessandri, A.\ et al.\ (2017). Multi-scale enhancement of climate prediction over land by increasing the model
sensitivity to vegetation variability in EC-Earth. \emph{Clim.\ Dyn.}\ 49, 1215--1237.\par
Beljaars, A.C.M.\ \& Viterbo, P.\ (1998). Role of the boundary layer in a numerical weather prediction model. In
\emph{Clear and Cloudy Boundary Layers}, KNAW, 287--304.\par
Betts, R.A.\ (2000). Offset of the potential carbon sink from boreal forestation by decreases in surface albedo.
\emph{Nature} 408, 187--190.\par
Bonan, G.B.\ (2008). Forests and climate change: forcings, feedbacks, and the climate benefits of forests.
\emph{Science} 320, 1444--1449.\par
Bonan, G.B., Pollard, D.\ \& Thompson, S.L.\ (1992). Effects of boreal forest vegetation on global climate.
\emph{Nature} 359, 716--718.\par
Chapin, F.S.\ et al.\ (2005). Role of land-surface changes in Arctic summer warming. \emph{Science} 310, 657--660.\par
de Lavergne, C.\ et al.\ (2014). Cessation of deep convection in the open Southern Ocean under anthropogenic climate
change. \emph{Nat.\ Clim.\ Change} 4, 278--282.\par
D\"oscher, R.\ et al.\ (2022). The EC-Earth3 Earth system model for CMIP6. \emph{Geosci.\ Model Dev.}\ 15,
2973--3020.\par
Herzschuh, U.\ (2020). Legacy of the Last Glacial on the present-day distribution of deciduous versus evergreen
boreal forests. \emph{Glob.\ Ecol.\ Biogeogr.}\ 29, 198--206.\par
Heuz\'e, C.\ (2021). Antarctic bottom water and North Atlantic deep water in CMIP6 models. \emph{Ocean Sci.}\ 17,
59--90.\par
Heuz\'e, C.\ et al.\ (2013). Southern Ocean bottom water characteristics in CMIP5 models. \emph{Geophys.\ Res.\
Lett.}\ 40, 1409--1414.\par
Holtslag, A.A.M.\ et al.\ (2013). Stable atmospheric boundary layers and diurnal cycles: challenges for weather and
climate models. \emph{Bull.\ Am.\ Meteorol.\ Soc.}\ 94, 1691--1706.\par
Hyder, P.\ et al.\ (2018). Critical Southern Ocean climate model biases traced to atmospheric model cloud errors.
\emph{Nat.\ Commun.}\ 9, 3625.\par
Jiang, Y.\ et al.\ (2012). [Boreal tree/grass fraction spread across dynamic global vegetation models.]
\emph{Ecol.\ Evol.}\ (structural-uncertainty reference; details to verify).\par
K\"orner, C.\ \& Paulsen, J.\ (2004). A world-wide study of high altitude treeline temperatures. \emph{J.\
Biogeogr.}\ 31, 713--732.\par
Loranty, M.M.\ et al.\ (2014). Vegetation controls on northern high-latitude snow-albedo feedback. \emph{Glob.\
Change Biol.}\ 20, 594--606.\par
Louis, J.-F.\ (1979). A parametric model of vertical eddy fluxes in the atmosphere. \emph{Bound.-Layer Meteorol.}\
17, 187--202.\par
Medeiros, B.\ et al.\ (2011). Arctic inversion strength in climate models. \emph{J.\ Climate} 24, 4733--4740.\par
Mohrmann, M., Heuz\'e, C.\ \& Swart, S.\ (2021). Southern Ocean polynyas in CMIP6 models. \emph{The Cryosphere} 15,
4281--4313.\par
Pithan, F., Medeiros, B.\ \& Mauritsen, T.\ (2014). Mixed-phase clouds cause climate model biases in Arctic
wintertime temperature inversions. \emph{Clim.\ Dyn.}\ 43, 289--303.\par
Sandu, I.\ et al.\ (2013). Why is it so difficult to represent stably stratified conditions in numerical weather
prediction models? \emph{J.\ Adv.\ Model.\ Earth Syst.}\ 5, 117--133.\par
Sitch, S.\ et al.\ (2003). Evaluation of ecosystem dynamics, plant geography and terrestrial carbon cycling in the
LPJ dynamic global vegetation model. \emph{Glob.\ Change Biol.}\ 9, 161--185.\par
Smith, B.\ et al.\ (2001). Representation of vegetation dynamics in the modelling of terrestrial ecosystems.
\emph{Glob.\ Ecol.\ Biogeogr.}\ 10, 621--637.\par
Smith, B.\ et al.\ (2014). Implications of incorporating N cycling and N limitations on primary production in an
individual-based dynamic vegetation model. \emph{Biogeosciences} 11, 2027--2054.\par
Zhang, Y.\ et al.\ (2011). On the computation of planetary boundary-layer height / polar inversion strength in
climate models. \emph{J.\ Climate} 24, 5167--5186.\par
\vspace{2pt}\emph{Citations added in response to internal review (rev.\ 2026-07-06); bibliographic details for
Beljaars \& Viterbo 1998, Jiang et al.\ 2012, Medeiros et al.\ 2011 and Zhang et al.\ 2011 to be page-verified.}
}

\end{document}
